% additional_sections.tex
% Draft sections for IFAC MIM 2027 / IEEE CASE 2027 submission
% Two new sections: ISA-95 Architecture Mapping and MESA Function 12 (AI/ML)
% Author: Stephen Eacuello — DRAFT for refinement
%
% Usage: % additional_sections.tex
% Draft sections for IFAC MIM 2027 / IEEE CASE 2027 submission
% Two new sections: ISA-95 Architecture Mapping and MESA Function 12 (AI/ML)
% Author: Stephen Eacuello — DRAFT for refinement
%
% Usage: % additional_sections.tex
% Draft sections for IFAC MIM 2027 / IEEE CASE 2027 submission
% Two new sections: ISA-95 Architecture Mapping and MESA Function 12 (AI/ML)
% Author: Stephen Eacuello — DRAFT for refinement
%
% Usage: % additional_sections.tex
% Draft sections for IFAC MIM 2027 / IEEE CASE 2027 submission
% Two new sections: ISA-95 Architecture Mapping and MESA Function 12 (AI/ML)
% Author: Stephen Eacuello — DRAFT for refinement
%
% Usage: \input{additional_sections} in the main document, or copy sections
% into Eacuello_MES_Design_Prototype.tex at the desired location.
%
% Placeholder citation keys:
%   \cite{mesa1997}      — MESA-11 White Paper #06, 1997
%   \cite{isa95}         — ANSI/ISA-95.00.01-2010 Part 1
%   \cite{isa95part3}    — ANSI/ISA-95.00.03 Part 3 (Activity Models of MOM)
%   \cite{isa95part4}    — ANSI/ISA-95.00.04 Part 4 (Objects and Attributes for MOM)
%   \cite{b2mml}         — B2MML (Business to Manufacturing Markup Language) schema
%   \cite{shojaeinasab2022} — Intelligent MES systematic review
%   \cite{mesa2004}      — MESA Strategic Initiatives Model, 2004
%   \cite{liu2017}       — Isolation Forest (Liu et al., 2008/2012)
%   \cite{montgomery2019} — Montgomery, Introduction to SQC, 8th ed.


%======================================================================
\section{ISA-95 Architecture Mapping}\label{sec:isa95mapping}
%======================================================================

The ISA-95 standard (ANSI/ISA-95.00.01-2010) defines a hierarchical model for
enterprise--control system integration spanning five levels: Level~0 (physical
process), Level~1 (sensing and actuation), Level~2 (supervisory control),
Level~3 (manufacturing operations management), and Level~4 (business planning
and logistics)~\cite{isa95}.  While Section~\ref{sec:overview} introduced this
hierarchy in functional terms, this section provides a formal compliance mapping
between ISA-95 activity domains and the SodhiCable MES implementation modules,
demonstrating how a single Flask-based prototype can simulate the full
automation pyramid.

\subsection{Hierarchical Mapping}

Table~\ref{tab:isa95matrix} presents the ISA-95 compliance matrix, mapping each
level to its activity domain, the corresponding SodhiCable implementation
module(s), and key capabilities realized in the prototype.  Level~3 is further
decomposed by the 11 MESA functions as defined in ISA-95 Part~3's Activity
Models of Manufacturing Operations Management~\cite{isa95part3}.

\begin{table}[H]\centering\small
\caption{ISA-95 Compliance Matrix --- SodhiCable MES Architecture}
\label{tab:isa95matrix}
\begin{tabularx}{\textwidth}{@{} l L{2.4cm} L{3.0cm} L{5.5cm} @{}}
\toprule
\textbf{ISA-95 Level} & \textbf{Activity Domain} & \textbf{Implementation Module(s)} & \textbf{Key Capabilities} \\
\midrule
\multirow{2}{*}{Level 0} & Physical Process & DES Engine & Discrete event simulation of 8-stage factory; compound mixing, wire drawing, extrusion, braiding, cabling, jacketing, testing, packaging \\
\addlinespace[3pt]
\multirow{2}{*}{Level 1} & Sensing \& Actuation & OPC-UA Tag Simulator; \tfield{process\_data\_live}; \tfield{spark\_test\_log}; \tfield{energy\_readings} & 9{,}200+ simulated sensor readings (thermocouples, laser micrometers, pressure transducers, line speed encoders, spark testers); configurable scan interval (default 5\,s) \\
\addlinespace[3pt]
\multirow{2}{*}{Level 2} & Supervisory Control (SCADA/PLC) & \texttt{api\_scada.py}; PID simulation; PLC status endpoint & SCADA drill-down (plant $\to$ work center $\to$ sensor); PLC control-loop visualization (AUTO/MANUAL mode, setpoint tracking, interlock status); alarm management with severity ranking; energy monitoring per work center \\
\addlinespace[3pt]
Level 3 --- F1 & Resource Allocation & \texttt{api\_resources.py} & LP assignment of 26 WCs to WOs; capacity tracking; equipment status management \\
Level 3 --- F2 & Operations Scheduling & \texttt{api\_schedule.py} & SPT/EDD/MILP sequencing; changeover matrix; Gantt visualization \\
Level 3 --- F3 & Dispatching & \texttt{api\_dispatch.py} & Weighted priority dispatch; queue management; real-time job release \\
Level 3 --- F4 & Document Control & \texttt{api\_documents.py} & Recipe versioning (ISA-88 hierarchy); C\,of\,C generation; SOP management \\
Level 3 --- F5 & Data Collection & \texttt{api\_datacollection.py} & Real-time sensor ingestion; threshold alarming; data collection point configuration \\
Level 3 --- F6 & Labor Management & \texttt{api\_labor.py} & Skill matrix; certification tracking; IP shift rostering \\
Level 3 --- F7 & Quality Management & \texttt{api\_quality.py} & $\bar{X}$-$R$ charts; $C_{pk}$ computation; NCR$\to$CAPA workflow; SPC readings \\
Level 3 --- F8 & Process Management & \texttt{api\_process.py} & Recipe execution; CUSUM/EWMA drift detection; PID simulation; hold/release \\
Level 3 --- F9 & Maintenance Mgmt & \texttt{api\_maintenance.py} & MTBF-based PM scheduling; corrective maintenance logging; calibration tracking \\
Level 3 --- F10 & Product Tracking & \texttt{api\_genealogy.py} & Recursive CTE forward/backward trace; lot quarantine; splice zone detection \\
Level 3 --- F11 & Performance Analysis & \texttt{api\_performance.py} & OEE = $A \times P \times Q$; shift reporting; scrap Pareto; KPI dashboards \\
\addlinespace[3pt]
\multirow{2}{*}{Level 4} & Business Planning (ERP) & \texttt{api\_erp.py}; ERP Simulator & S\&OP planning; purchase orders with goods receipt; order-to-cash pipeline; demand forecasting (SES/DES); invoicing; financial P\&L; ISA-95 B2MML message exchange \\
\bottomrule
\end{tabularx}
\end{table}


\subsection{B2MML Message Integration}

The ISA-95 standard specifies information exchange between Level~3 and Level~4
through structured message types defined in the Business to Manufacturing Markup
Language (B2MML)~\cite{b2mml}.  The SodhiCable prototype simulates these
exchanges through the \tfield{isa95\_messages} table, which logs every
cross-boundary transaction with direction, message type, source system, target
system, payload summary, and processing status.  Table~\ref{tab:b2mml} lists
the B2MML message types implemented in the prototype.

\begin{table}[H]\centering\small
\caption{B2MML Message Types Simulated in the SodhiCable MES}
\label{tab:b2mml}
\begin{tabularx}{\textwidth}{@{} l l l L{5.0cm} @{}}
\toprule
\textbf{Direction} & \textbf{B2MML Message Type} & \textbf{Trigger Event} & \textbf{Payload Description} \\
\midrule
L4$\to$L3 & ProductionSchedule & S\&OP approval; WO release from SO & Approved production plan released to MES for scheduling and dispatch \\
L4$\to$L3 & MaterialRequirement & Manual PO creation & Purchase order for raw materials with quantity, supplier, and expected delivery \\
L3$\to$L4 & MaterialReceipt & Goods receipt processing & Goods receipt confirmation with accepted/rejected quantities and lot number \\
L3$\to$L4 & ProductionPerformance & Shift report submission & OEE metrics, output quantities, scrap totals, and downtime summary per shift \\
% AUTHOR NOTE: Add additional message types as implemented, e.g.:
% L4$\to$L3 & MaterialDefinition & Product/BOM update & Material master data synchronization \\
% L3$\to$L4 & QualityTestResult & Test completion & Spark test, dimensional, and DC resistance results \\
\bottomrule
\end{tabularx}
\end{table}


\subsection{Cross-Level Integration Points}

Three cross-level integration patterns merit particular attention, as they
demonstrate how the MES bridges automation and enterprise layers:

\begin{enumerate}[leftmargin=2em, itemsep=2pt]
\item \textbf{Traceability spanning L1--L4.}  A spark test failure at Level~1
(sensor detects pinhole in \tfield{spark\_test\_log}) propagates upward through
Level~2 (alarm triggered, PLC interlock status set to TRIPPED in the
\texttt{plc\_status} endpoint), Level~3 (NCR auto-created in \tfield{ncr}, lot
quarantined in \tfield{inventory}, forward trace initiated via
\tfield{lot\_tracking}), and Level~4 (yield KPI decremented in
\tfield{shift\_reports}, customer promise date re-evaluated in
\tfield{sales\_orders}).  This end-to-end chain is exercised in the integration
timeline (Section~\ref{sec:timeline}, Events~5--8).

\item \textbf{Recipe--sensor closed loop spanning L2--L3.}  Recipe setpoints
defined at Level~3 (\tfield{recipe\_parameters}: CV zone temperatures, PLCV
pressure, line speed) are pushed to Level~2 as PLC setpoints.  The SCADA PLC
status endpoint (\texttt{/api/scada/plc\_status/\{wc\_id\}}) compares actual
sensor readings against recipe setpoints and computes control error, output
percentage, and interlock status.  Process deviations detected by CUSUM or EWMA
at Level~3 feed back as corrective actions or hold/release decisions.

\item \textbf{AI as intelligence overlay spanning L2--L4.}  The AI/ML layer
(Section~\ref{sec:ai}) operates as a cross-cutting intelligence service
consuming data from all ISA-95 levels: Level~1/2 sensor data for anomaly
detection (z-score, Isolation Forest, Mahalanobis distance), Level~3 production
data for quality prediction (Gradient Boosted Stumps) and maintenance
forecasting (Weibull RUL), and Level~4 demand and financial data for
AI-driven recommendations.  This cross-level positioning distinguishes AI
from the traditional MESA functions, which operate primarily within Level~3.
% AUTHOR NOTE: Consider adding a diagram showing the AI overlay across levels.
\end{enumerate}


%======================================================================
\section{Toward MESA Function~12: An AI/ML Intelligence Layer}\label{sec:mesa12}
%======================================================================

\subsection{Motivation: The MESA-11 Model and the AI Gap}

The MESA International 11-function model, published in 1997~\cite{mesa1997},
defined the canonical functional decomposition of Manufacturing Execution
Systems that has guided MES implementations for nearly three decades.  The
eleven functions---Resource Allocation (F1), Operations Scheduling (F2),
Dispatching (F3), Document Control (F4), Data Collection (F5), Labor Management
(F6), Quality Management (F7), Process Management (F8), Maintenance Management
(F9), Product Tracking (F10), and Performance Analysis (F11)---were conceived in
an era when shop-floor computing consisted of SCADA/DCS systems with limited
storage, batch-mode SPC calculations, and manual analysis of printed control
charts.

The 1997 framework implicitly assumes that human domain experts perform the
analytical synthesis: an engineer reviews control charts (F7), correlates them
with process logs (F8), cross-references maintenance history (F9), and makes
decisions.  The framework provides no function for automated pattern recognition,
predictive modeling, multivariate anomaly detection, or machine-learned
optimization---capabilities that modern AI/ML methods can now deliver at
production-data scale and latency.

MESA International recognized the need for expansion, releasing the Strategic
Initiatives Model in 2004~\cite{mesa2004} and the Smart Manufacturing Model in
subsequent years.  ISA-95 itself was extended with Part~4~\cite{isa95part4},
which added activity models for Manufacturing Operations Management (MOM)
including production operations, quality operations, maintenance operations, and
inventory operations.  However, neither MESA nor ISA-95 has formalized
an AI/ML function within the MES functional hierarchy.

We propose that the original MESA-11 model requires a twelfth function:
\textbf{F12~--- AI/ML Intelligence Layer}, a cross-cutting analytical function
that consumes data produced by F1--F11 and generates predictive, prescriptive,
and diagnostic outputs that augment human decision-making.  The SodhiCable
prototype provides an existence proof through eight implemented AI/ML
capabilities.


\subsection{Implemented AI/ML Capabilities}

The SodhiCable prototype implements eight distinct AI/ML capabilities in pure
Python (no external ML libraries), demonstrating that meaningful intelligence
can be embedded directly in a lightweight MES without requiring a separate data
science platform.  Table~\ref{tab:f12capabilities} summarizes the eight
capabilities and their MESA function touchpoints.

\begin{table}[H]\centering\small
\caption{Proposed F12 Capabilities Implemented in the SodhiCable MES Prototype}
\label{tab:f12capabilities}
\begin{tabularx}{\textwidth}{@{} c L{3.0cm} L{3.5cm} l L{2.5cm} @{}}
\toprule
\textbf{\#} & \textbf{Capability} & \textbf{Method} & \textbf{Touches} & \textbf{Data Source} \\
\midrule
1 & Univariate anomaly detection & Z-score ($|z| > 2.5\sigma$) & F5, F8 & \tfield{process\_data\_live} \\
2 & Multivariate anomaly detection & Isolation Forest (100-tree ensemble) & F5, F8 & \tfield{process\_data\_live} \\
3 & Quality prediction & Gradient Boosted Stumps (50-round additive ensemble) & F7, F11 & \tfield{process\_deviations}, \tfield{scrap\_log}, \tfield{shift\_reports} \\
4 & Multivariate anomaly detection & Mahalanobis distance ($D^2 > \chi^2_{\alpha,p}$) & F5, F7, F8 & \tfield{process\_data\_live} \\
5 & Predictive maintenance & Weibull/exponential RUL estimation & F9 & \tfield{equipment}, \tfield{maintenance}, \tfield{downtime\_log} \\
6 & Statistical process monitoring & CUSUM drift \& EWMA smoothing & F7, F8 & \tfield{process\_deviations}, \tfield{spc\_readings} \\
7 & Prescriptive recommendations & Rules-based engine (7 rule classes) & F1--F11 & Cross-function aggregation \\
8 & Natural language query & Claude API $\to$ SQL generation $\to$ execution & F5, F11 & Full 78-table schema \\
\bottomrule
\end{tabularx}
\end{table}

\subsubsection{Z-Score Anomaly Detection}

The simplest analytical capability: for each (work center, parameter) group in
\tfield{process\_data\_live}, the system computes the sample mean $\bar{x}$ and
standard deviation $s$, then flags any reading where $|z_i| = |x_i - \bar{x}|/s
> 2.5$ as anomalous.  Readings with $|z| > 3.0$ are classified as Critical;
those in the range $2.5 < |z| \leq 3.0$ as Warning.  This extends F5 (data
collection) by adding statistical context to raw sensor values and F8 (process
management) by detecting deviations that may not cross fixed threshold alarms
but represent statistically unusual behavior.

\subsubsection{Isolation Forest}

The prototype implements a pure-Python Isolation Forest~\cite{liu2017}
consisting of 100 isolation trees with a maximum depth of 8 and a subsample
size of 256.  For each (work center, parameter) group, a three-dimensional
feature vector is constructed: $[\,z\text{-score},\;\Delta v / s,\;\Delta v_{-2}
/ s\,]$, where $\Delta v$ is the first difference and $\Delta v_{-2}$ is the
second-order lag difference, both normalized by the parameter's standard
deviation.  The anomaly score is computed as $s(x, n) = 2^{-E[h(x)]/c(n)}$,
where $E[h(x)]$ is the average path length across the 100 trees and $c(n)$ is
the expected path length of an unsuccessful search in a binary search tree of
$n$ items.  Readings with $s > 0.65$ are flagged (Critical if $s > 0.75$).

This method detects interaction effects that univariate z-score analysis
misses---for example, a zone temperature and line speed each within normal range
individually, but whose combination is anomalous.  The implementation requires
approximately 80 lines of Python with no external dependencies.

\subsubsection{Gradient Boosted Decision Stumps}

A 50-round gradient boosting ensemble predicts scrap rate per work center.  Each
round fits a single-split decision stump (one-level tree) to the current
residuals, selecting the feature and split point that minimize mean squared
error.  Four features are extracted per work center: (i)~average deviation
percentage from setpoints, (ii)~count of unresolved process deviations,
(iii)~average OEE, and (iv)~cumulative downtime minutes.  The learning rate is
0.3, and the model trains in-memory on each API call (stateless, reproducible).

Feature importance is computed as the cumulative absolute difference between
left and right split values across all 50 rounds, normalized to percentages.
The RMSE on the training set provides a diagnostic for model fit.  Work centers
with predicted scrap rates above 5\% are flagged as Critical, and those above
3\% as Warning.

\subsubsection{Mahalanobis Distance}

For multivariate anomaly detection that accounts for inter-parameter
correlations, the prototype computes the Mahalanobis distance $D^2$ for each
observation vector:
\[
D^2 = (\mathbf{x} - \boldsymbol{\mu})^\top \, \boldsymbol{\Sigma}^{-1} \, (\mathbf{x} - \boldsymbol{\mu})
\]
where $\boldsymbol{\mu}$ is the sample mean vector and $\boldsymbol{\Sigma}$ is
the sample covariance matrix, computed from all readings for a given work
center.  Observations exceeding the $\chi^2$ critical value at the 99\%
confidence level ($\chi^2_{0.01,2} = 9.21$ for $p=2$ features;
$\chi^2_{0.01,3} = 11.34$ for $p=3$) are flagged as anomalous.

The covariance matrix inversion is implemented analytically for $p = 2$ and
$p = 3$ using the adjugate/determinant method (no NumPy dependency).  For work
centers with more than three tracked parameters, the analysis uses the first
two principal features to maintain tractability in pure Python.

\subsubsection{Weibull/Exponential RUL Prediction}

Remaining Useful Life (RUL) estimation uses a two-model approach.  When
sufficient time-between-failure (TBF) data is available ($\geq 3$
observations), a Weibull distribution is fitted using the method of moments:
\[
F(t) = 1 - e^{-(t/\eta)^\beta}
\]
where the shape parameter $\beta$ is estimated from the coefficient of variation
of the TBF data ($\beta \approx 1/\text{CV}$, clamped to $[\,0.5, \infty)$)
and the scale parameter $\eta$ from the sample mean.  The conditional RUL is
estimated as $E[T - t \mid T > t] \approx E[T] - t$ for the current elapsed
time $t$.  For equipment with insufficient failure history, the model falls back
to the memoryless exponential distribution ($\beta = 1$) with a default
MTBF of 45 days.

The 7-day failure probability $P(\text{fail in } [t, t+168\text{\,hr}])$ is
computed from the Weibull survival function difference, enabling risk-ranked
prioritization of preventive maintenance scheduling.

\subsubsection{CUSUM and EWMA Statistical Monitoring}

Cumulative Sum (CUSUM) and Exponentially Weighted Moving Average (EWMA) charts
extend traditional Shewhart-type SPC (F7/F8) by detecting sustained small
shifts that $\bar{X}$-$R$ charts miss~\cite{montgomery2019}.  The prototype
stores detected CUSUM drifts in \tfield{process\_deviations} with
\tfield{detection\_method = `CUSUM'} and tracks resolution status.  The AI
recommendations engine (Capability~7) monitors unresolved CUSUM drifts and
generates corrective action suggestions referencing likely root causes (die
wear, material variation, calibration drift).

\subsubsection{AI-Driven Recommendations Engine}

A rules-based recommendation engine synthesizes data from across all MESA
functions to generate prioritized, actionable recommendations.  Seven rule
classes are implemented:

\begin{enumerate}[leftmargin=2em, itemsep=1pt]
\item High scrap rate ($>5\%$) at any work center (touches F5, F7, F8, F10)
\item Unresolved CUSUM drift on process parameters (touches F5, F7, F8)
\item Overdue preventive maintenance (touches F1, F9)
\item Low process capability $C_{pk} < 1.33$ (touches F7, F8, F11)
\item Expiring operator certifications within 30 days (touches F1, F6)
\item Declining OEE trend (7-day rolling average) (touches F8, F9, F11)
\item Accumulation of unresolved process deviations ($>3$) (touches F7, F8)
\end{enumerate}

Each recommendation includes a priority level (1 = highest), affected work
center(s), a natural-language explanation with quantitative evidence, and the
MESA function codes it touches.  This cross-function synthesis---combining
quality, maintenance, process, labor, and performance data into a single ranked
action list---is precisely the kind of analytical integration that the original
MESA-11 model left to human judgment.

\subsubsection{Natural Language Query Interface}

The most advanced F12 capability integrates the Anthropic Claude API to provide
a natural language $\to$ SQL $\to$ natural language query pipeline.  The system
constructs a schema-aware prompt containing the full 78-table database structure
and wire-and-cable domain knowledge (OEE targets, $C_{pk}$ thresholds, product
family codes, work center types), then uses Claude's tool-use capability to
iteratively generate and execute SQL queries against the MES database.  Only
\texttt{SELECT} queries are permitted (enforced server-side), and results are
limited to 50 rows.

This capability democratizes MES data access: an operator can ask ``Which lot
should I prioritize?'' and receive a ranked dispatch queue, or a manager can ask
``How much scrap this week?'' and receive an aggregated summary by cause
code---without writing SQL or navigating multiple dashboard screens.  The
agentic loop supports up to 8 iterations of query refinement, enabling complex
multi-step analyses.
% AUTHOR NOTE: Consider adding a concrete query example with input/output
% to illustrate the interaction pattern.


\subsection{F12 as a Cross-Cutting Function}

Unlike functions F1--F11, which each address a distinct operational domain, the
proposed F12 is inherently \textit{cross-cutting}: it consumes data from and
provides analytical outputs to multiple existing functions simultaneously.
Figure~\ref{fig:f12touchpoints} (conceptual; to be drawn) would show F12 at the
center of a hub-and-spoke diagram, with connections to:

\begin{itemize}[leftmargin=2em, itemsep=1pt]
\item \textbf{F5 --- Data Collection:} F12 consumes the sensor data that F5
collects, adding statistical and ML-based anomaly detection as a value-added
analytical layer atop raw data ingestion.
\item \textbf{F7 --- Quality Management:} F12 extends SPC with predictive
quality models (Gradient Boosted Stumps predicting scrap rate from process
features) and multivariate anomaly detection (Mahalanobis distance capturing
inter-parameter correlations that univariate control charts miss).
\item \textbf{F8 --- Process Management:} F12 augments CUSUM/EWMA drift
detection with Isolation Forest anomaly scoring and feeds recommendations back
to process engineers for corrective action.
\item \textbf{F9 --- Maintenance Management:} F12 replaces fixed-interval PM
scheduling with condition-based RUL estimation (Weibull model), enabling
risk-ranked maintenance prioritization.
\item \textbf{F11 --- Performance Analysis:} F12 detects OEE trends and
correlates declining performance with root causes across quality, maintenance,
and process domains---an analytical synthesis that F11's descriptive KPIs alone
cannot provide.
\end{itemize}

This cross-cutting nature is analogous to how ISA-95 Part~4~\cite{isa95part4}
extended the original Parts~1--3 by defining Manufacturing Operations Management
(MOM) activity models that cut across production, quality, maintenance, and
inventory operations.  Just as Part~4 recognized that operational activities do
not exist in functional silos, F12 recognizes that analytical intelligence must
span the full MES data landscape to be effective.
% AUTHOR NOTE: A formal argument for standardization would require mapping F12
% to ISA-95 Part 3 activity models and demonstrating that no existing activity
% adequately covers AI/ML analytics.  Consider drafting this for the conference
% paper.


\subsection{Toward Standardization}

The proposal for MESA Function~12 is not merely academic.  Shojaeinasab et
al.~\cite{shojaeinasab2022} identify five intelligence levels for MES, from
reactive (Level~1) to cognitive (Level~5), and note that most industrial MES
implementations remain at Levels~1--2 (reactive and proactive).  The SodhiCable
prototype operates at Level~3 (predictive) through its Weibull RUL,
Gradient Boosted quality prediction, and anomaly detection capabilities, with
elements of Level~4 (prescriptive) in the recommendations engine.

We suggest that formalizing F12 within the MESA framework would provide three
benefits:

\begin{enumerate}[leftmargin=2em, itemsep=2pt]
\item \textbf{Architectural clarity.}  Vendors and implementers would have a
defined functional boundary for AI/ML capabilities, preventing both
under-investment (treating AI as an afterthought) and scope creep (embedding ML
in every function without coherent architecture).

\item \textbf{Data contract standardization.}  An F12 specification would define
the minimum data interfaces required between the intelligence layer and
functions F1--F11, analogous to how ISA-95 Part~2 defines object models for the
L3--L4 boundary.  This would enable plug-and-play substitution of AI/ML engines
(e.g., replacing the pure-Python Isolation Forest with a scikit-learn or
cloud-based model) without modifying the MES core.

\item \textbf{Maturity assessment.}  An F12 specification could incorporate the
intelligence-level taxonomy of Shojaeinasab et al.~\cite{shojaeinasab2022},
providing a maturity model for organizations to assess and benchmark their MES
AI capabilities.
\end{enumerate}

The SodhiCable prototype demonstrates that a meaningful F12 implementation is
feasible with minimal infrastructure: all eight capabilities run in pure Python
within a single Flask application, with no external ML libraries, GPU
requirements, or separate data platforms.  This accessibility is critical for
the small and mid-size manufacturers (SMMs) that constitute the majority of the
wire and cable industry and for whom heavyweight AI infrastructure is not
economically viable.

% AUTHOR NOTE: For the conference submission, consider:
% (1) Adding empirical validation --- run the AI endpoints against the seeded
%     data and report precision/recall for anomaly detection, RMSE for quality
%     prediction, and RUL calibration accuracy.
% (2) Comparing F12 outputs against human expert decisions on the same data
%     to quantify the value of automated intelligence.
% (3) Surveying MESA International's current working groups to determine if
%     an AI/ML function is already under discussion.
% (4) Discussing computational complexity and scalability --- the pure-Python
%     implementations are O(n * n_trees * max_depth) for Isolation Forest;
%     for production scale, discuss when external ML infrastructure becomes
%     necessary.
 in the main document, or copy sections
% into Eacuello_MES_Design_Prototype.tex at the desired location.
%
% Placeholder citation keys:
%   \cite{mesa1997}      — MESA-11 White Paper #06, 1997
%   \cite{isa95}         — ANSI/ISA-95.00.01-2010 Part 1
%   \cite{isa95part3}    — ANSI/ISA-95.00.03 Part 3 (Activity Models of MOM)
%   \cite{isa95part4}    — ANSI/ISA-95.00.04 Part 4 (Objects and Attributes for MOM)
%   \cite{b2mml}         — B2MML (Business to Manufacturing Markup Language) schema
%   \cite{shojaeinasab2022} — Intelligent MES systematic review
%   \cite{mesa2004}      — MESA Strategic Initiatives Model, 2004
%   \cite{liu2017}       — Isolation Forest (Liu et al., 2008/2012)
%   \cite{montgomery2019} — Montgomery, Introduction to SQC, 8th ed.


%======================================================================
\section{ISA-95 Architecture Mapping}\label{sec:isa95mapping}
%======================================================================

The ISA-95 standard (ANSI/ISA-95.00.01-2010) defines a hierarchical model for
enterprise--control system integration spanning five levels: Level~0 (physical
process), Level~1 (sensing and actuation), Level~2 (supervisory control),
Level~3 (manufacturing operations management), and Level~4 (business planning
and logistics)~\cite{isa95}.  While Section~\ref{sec:overview} introduced this
hierarchy in functional terms, this section provides a formal compliance mapping
between ISA-95 activity domains and the SodhiCable MES implementation modules,
demonstrating how a single Flask-based prototype can simulate the full
automation pyramid.

\subsection{Hierarchical Mapping}

Table~\ref{tab:isa95matrix} presents the ISA-95 compliance matrix, mapping each
level to its activity domain, the corresponding SodhiCable implementation
module(s), and key capabilities realized in the prototype.  Level~3 is further
decomposed by the 11 MESA functions as defined in ISA-95 Part~3's Activity
Models of Manufacturing Operations Management~\cite{isa95part3}.

\begin{table}[H]\centering\small
\caption{ISA-95 Compliance Matrix --- SodhiCable MES Architecture}
\label{tab:isa95matrix}
\begin{tabularx}{\textwidth}{@{} l L{2.4cm} L{3.0cm} L{5.5cm} @{}}
\toprule
\textbf{ISA-95 Level} & \textbf{Activity Domain} & \textbf{Implementation Module(s)} & \textbf{Key Capabilities} \\
\midrule
\multirow{2}{*}{Level 0} & Physical Process & DES Engine & Discrete event simulation of 8-stage factory; compound mixing, wire drawing, extrusion, braiding, cabling, jacketing, testing, packaging \\
\addlinespace[3pt]
\multirow{2}{*}{Level 1} & Sensing \& Actuation & OPC-UA Tag Simulator; \tfield{process\_data\_live}; \tfield{spark\_test\_log}; \tfield{energy\_readings} & 9{,}200+ simulated sensor readings (thermocouples, laser micrometers, pressure transducers, line speed encoders, spark testers); configurable scan interval (default 5\,s) \\
\addlinespace[3pt]
\multirow{2}{*}{Level 2} & Supervisory Control (SCADA/PLC) & \texttt{api\_scada.py}; PID simulation; PLC status endpoint & SCADA drill-down (plant $\to$ work center $\to$ sensor); PLC control-loop visualization (AUTO/MANUAL mode, setpoint tracking, interlock status); alarm management with severity ranking; energy monitoring per work center \\
\addlinespace[3pt]
Level 3 --- F1 & Resource Allocation & \texttt{api\_resources.py} & LP assignment of 26 WCs to WOs; capacity tracking; equipment status management \\
Level 3 --- F2 & Operations Scheduling & \texttt{api\_schedule.py} & SPT/EDD/MILP sequencing; changeover matrix; Gantt visualization \\
Level 3 --- F3 & Dispatching & \texttt{api\_dispatch.py} & Weighted priority dispatch; queue management; real-time job release \\
Level 3 --- F4 & Document Control & \texttt{api\_documents.py} & Recipe versioning (ISA-88 hierarchy); C\,of\,C generation; SOP management \\
Level 3 --- F5 & Data Collection & \texttt{api\_datacollection.py} & Real-time sensor ingestion; threshold alarming; data collection point configuration \\
Level 3 --- F6 & Labor Management & \texttt{api\_labor.py} & Skill matrix; certification tracking; IP shift rostering \\
Level 3 --- F7 & Quality Management & \texttt{api\_quality.py} & $\bar{X}$-$R$ charts; $C_{pk}$ computation; NCR$\to$CAPA workflow; SPC readings \\
Level 3 --- F8 & Process Management & \texttt{api\_process.py} & Recipe execution; CUSUM/EWMA drift detection; PID simulation; hold/release \\
Level 3 --- F9 & Maintenance Mgmt & \texttt{api\_maintenance.py} & MTBF-based PM scheduling; corrective maintenance logging; calibration tracking \\
Level 3 --- F10 & Product Tracking & \texttt{api\_genealogy.py} & Recursive CTE forward/backward trace; lot quarantine; splice zone detection \\
Level 3 --- F11 & Performance Analysis & \texttt{api\_performance.py} & OEE = $A \times P \times Q$; shift reporting; scrap Pareto; KPI dashboards \\
\addlinespace[3pt]
\multirow{2}{*}{Level 4} & Business Planning (ERP) & \texttt{api\_erp.py}; ERP Simulator & S\&OP planning; purchase orders with goods receipt; order-to-cash pipeline; demand forecasting (SES/DES); invoicing; financial P\&L; ISA-95 B2MML message exchange \\
\bottomrule
\end{tabularx}
\end{table}


\subsection{B2MML Message Integration}

The ISA-95 standard specifies information exchange between Level~3 and Level~4
through structured message types defined in the Business to Manufacturing Markup
Language (B2MML)~\cite{b2mml}.  The SodhiCable prototype simulates these
exchanges through the \tfield{isa95\_messages} table, which logs every
cross-boundary transaction with direction, message type, source system, target
system, payload summary, and processing status.  Table~\ref{tab:b2mml} lists
the B2MML message types implemented in the prototype.

\begin{table}[H]\centering\small
\caption{B2MML Message Types Simulated in the SodhiCable MES}
\label{tab:b2mml}
\begin{tabularx}{\textwidth}{@{} l l l L{5.0cm} @{}}
\toprule
\textbf{Direction} & \textbf{B2MML Message Type} & \textbf{Trigger Event} & \textbf{Payload Description} \\
\midrule
L4$\to$L3 & ProductionSchedule & S\&OP approval; WO release from SO & Approved production plan released to MES for scheduling and dispatch \\
L4$\to$L3 & MaterialRequirement & Manual PO creation & Purchase order for raw materials with quantity, supplier, and expected delivery \\
L3$\to$L4 & MaterialReceipt & Goods receipt processing & Goods receipt confirmation with accepted/rejected quantities and lot number \\
L3$\to$L4 & ProductionPerformance & Shift report submission & OEE metrics, output quantities, scrap totals, and downtime summary per shift \\
% AUTHOR NOTE: Add additional message types as implemented, e.g.:
% L4$\to$L3 & MaterialDefinition & Product/BOM update & Material master data synchronization \\
% L3$\to$L4 & QualityTestResult & Test completion & Spark test, dimensional, and DC resistance results \\
\bottomrule
\end{tabularx}
\end{table}


\subsection{Cross-Level Integration Points}

Three cross-level integration patterns merit particular attention, as they
demonstrate how the MES bridges automation and enterprise layers:

\begin{enumerate}[leftmargin=2em, itemsep=2pt]
\item \textbf{Traceability spanning L1--L4.}  A spark test failure at Level~1
(sensor detects pinhole in \tfield{spark\_test\_log}) propagates upward through
Level~2 (alarm triggered, PLC interlock status set to TRIPPED in the
\texttt{plc\_status} endpoint), Level~3 (NCR auto-created in \tfield{ncr}, lot
quarantined in \tfield{inventory}, forward trace initiated via
\tfield{lot\_tracking}), and Level~4 (yield KPI decremented in
\tfield{shift\_reports}, customer promise date re-evaluated in
\tfield{sales\_orders}).  This end-to-end chain is exercised in the integration
timeline (Section~\ref{sec:timeline}, Events~5--8).

\item \textbf{Recipe--sensor closed loop spanning L2--L3.}  Recipe setpoints
defined at Level~3 (\tfield{recipe\_parameters}: CV zone temperatures, PLCV
pressure, line speed) are pushed to Level~2 as PLC setpoints.  The SCADA PLC
status endpoint (\texttt{/api/scada/plc\_status/\{wc\_id\}}) compares actual
sensor readings against recipe setpoints and computes control error, output
percentage, and interlock status.  Process deviations detected by CUSUM or EWMA
at Level~3 feed back as corrective actions or hold/release decisions.

\item \textbf{AI as intelligence overlay spanning L2--L4.}  The AI/ML layer
(Section~\ref{sec:ai}) operates as a cross-cutting intelligence service
consuming data from all ISA-95 levels: Level~1/2 sensor data for anomaly
detection (z-score, Isolation Forest, Mahalanobis distance), Level~3 production
data for quality prediction (Gradient Boosted Stumps) and maintenance
forecasting (Weibull RUL), and Level~4 demand and financial data for
AI-driven recommendations.  This cross-level positioning distinguishes AI
from the traditional MESA functions, which operate primarily within Level~3.
% AUTHOR NOTE: Consider adding a diagram showing the AI overlay across levels.
\end{enumerate}


%======================================================================
\section{Toward MESA Function~12: An AI/ML Intelligence Layer}\label{sec:mesa12}
%======================================================================

\subsection{Motivation: The MESA-11 Model and the AI Gap}

The MESA International 11-function model, published in 1997~\cite{mesa1997},
defined the canonical functional decomposition of Manufacturing Execution
Systems that has guided MES implementations for nearly three decades.  The
eleven functions---Resource Allocation (F1), Operations Scheduling (F2),
Dispatching (F3), Document Control (F4), Data Collection (F5), Labor Management
(F6), Quality Management (F7), Process Management (F8), Maintenance Management
(F9), Product Tracking (F10), and Performance Analysis (F11)---were conceived in
an era when shop-floor computing consisted of SCADA/DCS systems with limited
storage, batch-mode SPC calculations, and manual analysis of printed control
charts.

The 1997 framework implicitly assumes that human domain experts perform the
analytical synthesis: an engineer reviews control charts (F7), correlates them
with process logs (F8), cross-references maintenance history (F9), and makes
decisions.  The framework provides no function for automated pattern recognition,
predictive modeling, multivariate anomaly detection, or machine-learned
optimization---capabilities that modern AI/ML methods can now deliver at
production-data scale and latency.

MESA International recognized the need for expansion, releasing the Strategic
Initiatives Model in 2004~\cite{mesa2004} and the Smart Manufacturing Model in
subsequent years.  ISA-95 itself was extended with Part~4~\cite{isa95part4},
which added activity models for Manufacturing Operations Management (MOM)
including production operations, quality operations, maintenance operations, and
inventory operations.  However, neither MESA nor ISA-95 has formalized
an AI/ML function within the MES functional hierarchy.

We propose that the original MESA-11 model requires a twelfth function:
\textbf{F12~--- AI/ML Intelligence Layer}, a cross-cutting analytical function
that consumes data produced by F1--F11 and generates predictive, prescriptive,
and diagnostic outputs that augment human decision-making.  The SodhiCable
prototype provides an existence proof through eight implemented AI/ML
capabilities.


\subsection{Implemented AI/ML Capabilities}

The SodhiCable prototype implements eight distinct AI/ML capabilities in pure
Python (no external ML libraries), demonstrating that meaningful intelligence
can be embedded directly in a lightweight MES without requiring a separate data
science platform.  Table~\ref{tab:f12capabilities} summarizes the eight
capabilities and their MESA function touchpoints.

\begin{table}[H]\centering\small
\caption{Proposed F12 Capabilities Implemented in the SodhiCable MES Prototype}
\label{tab:f12capabilities}
\begin{tabularx}{\textwidth}{@{} c L{3.0cm} L{3.5cm} l L{2.5cm} @{}}
\toprule
\textbf{\#} & \textbf{Capability} & \textbf{Method} & \textbf{Touches} & \textbf{Data Source} \\
\midrule
1 & Univariate anomaly detection & Z-score ($|z| > 2.5\sigma$) & F5, F8 & \tfield{process\_data\_live} \\
2 & Multivariate anomaly detection & Isolation Forest (100-tree ensemble) & F5, F8 & \tfield{process\_data\_live} \\
3 & Quality prediction & Gradient Boosted Stumps (50-round additive ensemble) & F7, F11 & \tfield{process\_deviations}, \tfield{scrap\_log}, \tfield{shift\_reports} \\
4 & Multivariate anomaly detection & Mahalanobis distance ($D^2 > \chi^2_{\alpha,p}$) & F5, F7, F8 & \tfield{process\_data\_live} \\
5 & Predictive maintenance & Weibull/exponential RUL estimation & F9 & \tfield{equipment}, \tfield{maintenance}, \tfield{downtime\_log} \\
6 & Statistical process monitoring & CUSUM drift \& EWMA smoothing & F7, F8 & \tfield{process\_deviations}, \tfield{spc\_readings} \\
7 & Prescriptive recommendations & Rules-based engine (7 rule classes) & F1--F11 & Cross-function aggregation \\
8 & Natural language query & Claude API $\to$ SQL generation $\to$ execution & F5, F11 & Full 78-table schema \\
\bottomrule
\end{tabularx}
\end{table}

\subsubsection{Z-Score Anomaly Detection}

The simplest analytical capability: for each (work center, parameter) group in
\tfield{process\_data\_live}, the system computes the sample mean $\bar{x}$ and
standard deviation $s$, then flags any reading where $|z_i| = |x_i - \bar{x}|/s
> 2.5$ as anomalous.  Readings with $|z| > 3.0$ are classified as Critical;
those in the range $2.5 < |z| \leq 3.0$ as Warning.  This extends F5 (data
collection) by adding statistical context to raw sensor values and F8 (process
management) by detecting deviations that may not cross fixed threshold alarms
but represent statistically unusual behavior.

\subsubsection{Isolation Forest}

The prototype implements a pure-Python Isolation Forest~\cite{liu2017}
consisting of 100 isolation trees with a maximum depth of 8 and a subsample
size of 256.  For each (work center, parameter) group, a three-dimensional
feature vector is constructed: $[\,z\text{-score},\;\Delta v / s,\;\Delta v_{-2}
/ s\,]$, where $\Delta v$ is the first difference and $\Delta v_{-2}$ is the
second-order lag difference, both normalized by the parameter's standard
deviation.  The anomaly score is computed as $s(x, n) = 2^{-E[h(x)]/c(n)}$,
where $E[h(x)]$ is the average path length across the 100 trees and $c(n)$ is
the expected path length of an unsuccessful search in a binary search tree of
$n$ items.  Readings with $s > 0.65$ are flagged (Critical if $s > 0.75$).

This method detects interaction effects that univariate z-score analysis
misses---for example, a zone temperature and line speed each within normal range
individually, but whose combination is anomalous.  The implementation requires
approximately 80 lines of Python with no external dependencies.

\subsubsection{Gradient Boosted Decision Stumps}

A 50-round gradient boosting ensemble predicts scrap rate per work center.  Each
round fits a single-split decision stump (one-level tree) to the current
residuals, selecting the feature and split point that minimize mean squared
error.  Four features are extracted per work center: (i)~average deviation
percentage from setpoints, (ii)~count of unresolved process deviations,
(iii)~average OEE, and (iv)~cumulative downtime minutes.  The learning rate is
0.3, and the model trains in-memory on each API call (stateless, reproducible).

Feature importance is computed as the cumulative absolute difference between
left and right split values across all 50 rounds, normalized to percentages.
The RMSE on the training set provides a diagnostic for model fit.  Work centers
with predicted scrap rates above 5\% are flagged as Critical, and those above
3\% as Warning.

\subsubsection{Mahalanobis Distance}

For multivariate anomaly detection that accounts for inter-parameter
correlations, the prototype computes the Mahalanobis distance $D^2$ for each
observation vector:
\[
D^2 = (\mathbf{x} - \boldsymbol{\mu})^\top \, \boldsymbol{\Sigma}^{-1} \, (\mathbf{x} - \boldsymbol{\mu})
\]
where $\boldsymbol{\mu}$ is the sample mean vector and $\boldsymbol{\Sigma}$ is
the sample covariance matrix, computed from all readings for a given work
center.  Observations exceeding the $\chi^2$ critical value at the 99\%
confidence level ($\chi^2_{0.01,2} = 9.21$ for $p=2$ features;
$\chi^2_{0.01,3} = 11.34$ for $p=3$) are flagged as anomalous.

The covariance matrix inversion is implemented analytically for $p = 2$ and
$p = 3$ using the adjugate/determinant method (no NumPy dependency).  For work
centers with more than three tracked parameters, the analysis uses the first
two principal features to maintain tractability in pure Python.

\subsubsection{Weibull/Exponential RUL Prediction}

Remaining Useful Life (RUL) estimation uses a two-model approach.  When
sufficient time-between-failure (TBF) data is available ($\geq 3$
observations), a Weibull distribution is fitted using the method of moments:
\[
F(t) = 1 - e^{-(t/\eta)^\beta}
\]
where the shape parameter $\beta$ is estimated from the coefficient of variation
of the TBF data ($\beta \approx 1/\text{CV}$, clamped to $[\,0.5, \infty)$)
and the scale parameter $\eta$ from the sample mean.  The conditional RUL is
estimated as $E[T - t \mid T > t] \approx E[T] - t$ for the current elapsed
time $t$.  For equipment with insufficient failure history, the model falls back
to the memoryless exponential distribution ($\beta = 1$) with a default
MTBF of 45 days.

The 7-day failure probability $P(\text{fail in } [t, t+168\text{\,hr}])$ is
computed from the Weibull survival function difference, enabling risk-ranked
prioritization of preventive maintenance scheduling.

\subsubsection{CUSUM and EWMA Statistical Monitoring}

Cumulative Sum (CUSUM) and Exponentially Weighted Moving Average (EWMA) charts
extend traditional Shewhart-type SPC (F7/F8) by detecting sustained small
shifts that $\bar{X}$-$R$ charts miss~\cite{montgomery2019}.  The prototype
stores detected CUSUM drifts in \tfield{process\_deviations} with
\tfield{detection\_method = `CUSUM'} and tracks resolution status.  The AI
recommendations engine (Capability~7) monitors unresolved CUSUM drifts and
generates corrective action suggestions referencing likely root causes (die
wear, material variation, calibration drift).

\subsubsection{AI-Driven Recommendations Engine}

A rules-based recommendation engine synthesizes data from across all MESA
functions to generate prioritized, actionable recommendations.  Seven rule
classes are implemented:

\begin{enumerate}[leftmargin=2em, itemsep=1pt]
\item High scrap rate ($>5\%$) at any work center (touches F5, F7, F8, F10)
\item Unresolved CUSUM drift on process parameters (touches F5, F7, F8)
\item Overdue preventive maintenance (touches F1, F9)
\item Low process capability $C_{pk} < 1.33$ (touches F7, F8, F11)
\item Expiring operator certifications within 30 days (touches F1, F6)
\item Declining OEE trend (7-day rolling average) (touches F8, F9, F11)
\item Accumulation of unresolved process deviations ($>3$) (touches F7, F8)
\end{enumerate}

Each recommendation includes a priority level (1 = highest), affected work
center(s), a natural-language explanation with quantitative evidence, and the
MESA function codes it touches.  This cross-function synthesis---combining
quality, maintenance, process, labor, and performance data into a single ranked
action list---is precisely the kind of analytical integration that the original
MESA-11 model left to human judgment.

\subsubsection{Natural Language Query Interface}

The most advanced F12 capability integrates the Anthropic Claude API to provide
a natural language $\to$ SQL $\to$ natural language query pipeline.  The system
constructs a schema-aware prompt containing the full 78-table database structure
and wire-and-cable domain knowledge (OEE targets, $C_{pk}$ thresholds, product
family codes, work center types), then uses Claude's tool-use capability to
iteratively generate and execute SQL queries against the MES database.  Only
\texttt{SELECT} queries are permitted (enforced server-side), and results are
limited to 50 rows.

This capability democratizes MES data access: an operator can ask ``Which lot
should I prioritize?'' and receive a ranked dispatch queue, or a manager can ask
``How much scrap this week?'' and receive an aggregated summary by cause
code---without writing SQL or navigating multiple dashboard screens.  The
agentic loop supports up to 8 iterations of query refinement, enabling complex
multi-step analyses.
% AUTHOR NOTE: Consider adding a concrete query example with input/output
% to illustrate the interaction pattern.


\subsection{F12 as a Cross-Cutting Function}

Unlike functions F1--F11, which each address a distinct operational domain, the
proposed F12 is inherently \textit{cross-cutting}: it consumes data from and
provides analytical outputs to multiple existing functions simultaneously.
Figure~\ref{fig:f12touchpoints} (conceptual; to be drawn) would show F12 at the
center of a hub-and-spoke diagram, with connections to:

\begin{itemize}[leftmargin=2em, itemsep=1pt]
\item \textbf{F5 --- Data Collection:} F12 consumes the sensor data that F5
collects, adding statistical and ML-based anomaly detection as a value-added
analytical layer atop raw data ingestion.
\item \textbf{F7 --- Quality Management:} F12 extends SPC with predictive
quality models (Gradient Boosted Stumps predicting scrap rate from process
features) and multivariate anomaly detection (Mahalanobis distance capturing
inter-parameter correlations that univariate control charts miss).
\item \textbf{F8 --- Process Management:} F12 augments CUSUM/EWMA drift
detection with Isolation Forest anomaly scoring and feeds recommendations back
to process engineers for corrective action.
\item \textbf{F9 --- Maintenance Management:} F12 replaces fixed-interval PM
scheduling with condition-based RUL estimation (Weibull model), enabling
risk-ranked maintenance prioritization.
\item \textbf{F11 --- Performance Analysis:} F12 detects OEE trends and
correlates declining performance with root causes across quality, maintenance,
and process domains---an analytical synthesis that F11's descriptive KPIs alone
cannot provide.
\end{itemize}

This cross-cutting nature is analogous to how ISA-95 Part~4~\cite{isa95part4}
extended the original Parts~1--3 by defining Manufacturing Operations Management
(MOM) activity models that cut across production, quality, maintenance, and
inventory operations.  Just as Part~4 recognized that operational activities do
not exist in functional silos, F12 recognizes that analytical intelligence must
span the full MES data landscape to be effective.
% AUTHOR NOTE: A formal argument for standardization would require mapping F12
% to ISA-95 Part 3 activity models and demonstrating that no existing activity
% adequately covers AI/ML analytics.  Consider drafting this for the conference
% paper.


\subsection{Toward Standardization}

The proposal for MESA Function~12 is not merely academic.  Shojaeinasab et
al.~\cite{shojaeinasab2022} identify five intelligence levels for MES, from
reactive (Level~1) to cognitive (Level~5), and note that most industrial MES
implementations remain at Levels~1--2 (reactive and proactive).  The SodhiCable
prototype operates at Level~3 (predictive) through its Weibull RUL,
Gradient Boosted quality prediction, and anomaly detection capabilities, with
elements of Level~4 (prescriptive) in the recommendations engine.

We suggest that formalizing F12 within the MESA framework would provide three
benefits:

\begin{enumerate}[leftmargin=2em, itemsep=2pt]
\item \textbf{Architectural clarity.}  Vendors and implementers would have a
defined functional boundary for AI/ML capabilities, preventing both
under-investment (treating AI as an afterthought) and scope creep (embedding ML
in every function without coherent architecture).

\item \textbf{Data contract standardization.}  An F12 specification would define
the minimum data interfaces required between the intelligence layer and
functions F1--F11, analogous to how ISA-95 Part~2 defines object models for the
L3--L4 boundary.  This would enable plug-and-play substitution of AI/ML engines
(e.g., replacing the pure-Python Isolation Forest with a scikit-learn or
cloud-based model) without modifying the MES core.

\item \textbf{Maturity assessment.}  An F12 specification could incorporate the
intelligence-level taxonomy of Shojaeinasab et al.~\cite{shojaeinasab2022},
providing a maturity model for organizations to assess and benchmark their MES
AI capabilities.
\end{enumerate}

The SodhiCable prototype demonstrates that a meaningful F12 implementation is
feasible with minimal infrastructure: all eight capabilities run in pure Python
within a single Flask application, with no external ML libraries, GPU
requirements, or separate data platforms.  This accessibility is critical for
the small and mid-size manufacturers (SMMs) that constitute the majority of the
wire and cable industry and for whom heavyweight AI infrastructure is not
economically viable.

% AUTHOR NOTE: For the conference submission, consider:
% (1) Adding empirical validation --- run the AI endpoints against the seeded
%     data and report precision/recall for anomaly detection, RMSE for quality
%     prediction, and RUL calibration accuracy.
% (2) Comparing F12 outputs against human expert decisions on the same data
%     to quantify the value of automated intelligence.
% (3) Surveying MESA International's current working groups to determine if
%     an AI/ML function is already under discussion.
% (4) Discussing computational complexity and scalability --- the pure-Python
%     implementations are O(n * n_trees * max_depth) for Isolation Forest;
%     for production scale, discuss when external ML infrastructure becomes
%     necessary.
 in the main document, or copy sections
% into Eacuello_MES_Design_Prototype.tex at the desired location.
%
% Placeholder citation keys:
%   \cite{mesa1997}      — MESA-11 White Paper #06, 1997
%   \cite{isa95}         — ANSI/ISA-95.00.01-2010 Part 1
%   \cite{isa95part3}    — ANSI/ISA-95.00.03 Part 3 (Activity Models of MOM)
%   \cite{isa95part4}    — ANSI/ISA-95.00.04 Part 4 (Objects and Attributes for MOM)
%   \cite{b2mml}         — B2MML (Business to Manufacturing Markup Language) schema
%   \cite{shojaeinasab2022} — Intelligent MES systematic review
%   \cite{mesa2004}      — MESA Strategic Initiatives Model, 2004
%   \cite{liu2017}       — Isolation Forest (Liu et al., 2008/2012)
%   \cite{montgomery2019} — Montgomery, Introduction to SQC, 8th ed.


%======================================================================
\section{ISA-95 Architecture Mapping}\label{sec:isa95mapping}
%======================================================================

The ISA-95 standard (ANSI/ISA-95.00.01-2010) defines a hierarchical model for
enterprise--control system integration spanning five levels: Level~0 (physical
process), Level~1 (sensing and actuation), Level~2 (supervisory control),
Level~3 (manufacturing operations management), and Level~4 (business planning
and logistics)~\cite{isa95}.  While Section~\ref{sec:overview} introduced this
hierarchy in functional terms, this section provides a formal compliance mapping
between ISA-95 activity domains and the SodhiCable MES implementation modules,
demonstrating how a single Flask-based prototype can simulate the full
automation pyramid.

\subsection{Hierarchical Mapping}

Table~\ref{tab:isa95matrix} presents the ISA-95 compliance matrix, mapping each
level to its activity domain, the corresponding SodhiCable implementation
module(s), and key capabilities realized in the prototype.  Level~3 is further
decomposed by the 11 MESA functions as defined in ISA-95 Part~3's Activity
Models of Manufacturing Operations Management~\cite{isa95part3}.

\begin{table}[H]\centering\small
\caption{ISA-95 Compliance Matrix --- SodhiCable MES Architecture}
\label{tab:isa95matrix}
\begin{tabularx}{\textwidth}{@{} l L{2.4cm} L{3.0cm} L{5.5cm} @{}}
\toprule
\textbf{ISA-95 Level} & \textbf{Activity Domain} & \textbf{Implementation Module(s)} & \textbf{Key Capabilities} \\
\midrule
\multirow{2}{*}{Level 0} & Physical Process & DES Engine & Discrete event simulation of 8-stage factory; compound mixing, wire drawing, extrusion, braiding, cabling, jacketing, testing, packaging \\
\addlinespace[3pt]
\multirow{2}{*}{Level 1} & Sensing \& Actuation & OPC-UA Tag Simulator; \tfield{process\_data\_live}; \tfield{spark\_test\_log}; \tfield{energy\_readings} & 9{,}200+ simulated sensor readings (thermocouples, laser micrometers, pressure transducers, line speed encoders, spark testers); configurable scan interval (default 5\,s) \\
\addlinespace[3pt]
\multirow{2}{*}{Level 2} & Supervisory Control (SCADA/PLC) & \texttt{api\_scada.py}; PID simulation; PLC status endpoint & SCADA drill-down (plant $\to$ work center $\to$ sensor); PLC control-loop visualization (AUTO/MANUAL mode, setpoint tracking, interlock status); alarm management with severity ranking; energy monitoring per work center \\
\addlinespace[3pt]
Level 3 --- F1 & Resource Allocation & \texttt{api\_resources.py} & LP assignment of 26 WCs to WOs; capacity tracking; equipment status management \\
Level 3 --- F2 & Operations Scheduling & \texttt{api\_schedule.py} & SPT/EDD/MILP sequencing; changeover matrix; Gantt visualization \\
Level 3 --- F3 & Dispatching & \texttt{api\_dispatch.py} & Weighted priority dispatch; queue management; real-time job release \\
Level 3 --- F4 & Document Control & \texttt{api\_documents.py} & Recipe versioning (ISA-88 hierarchy); C\,of\,C generation; SOP management \\
Level 3 --- F5 & Data Collection & \texttt{api\_datacollection.py} & Real-time sensor ingestion; threshold alarming; data collection point configuration \\
Level 3 --- F6 & Labor Management & \texttt{api\_labor.py} & Skill matrix; certification tracking; IP shift rostering \\
Level 3 --- F7 & Quality Management & \texttt{api\_quality.py} & $\bar{X}$-$R$ charts; $C_{pk}$ computation; NCR$\to$CAPA workflow; SPC readings \\
Level 3 --- F8 & Process Management & \texttt{api\_process.py} & Recipe execution; CUSUM/EWMA drift detection; PID simulation; hold/release \\
Level 3 --- F9 & Maintenance Mgmt & \texttt{api\_maintenance.py} & MTBF-based PM scheduling; corrective maintenance logging; calibration tracking \\
Level 3 --- F10 & Product Tracking & \texttt{api\_genealogy.py} & Recursive CTE forward/backward trace; lot quarantine; splice zone detection \\
Level 3 --- F11 & Performance Analysis & \texttt{api\_performance.py} & OEE = $A \times P \times Q$; shift reporting; scrap Pareto; KPI dashboards \\
\addlinespace[3pt]
\multirow{2}{*}{Level 4} & Business Planning (ERP) & \texttt{api\_erp.py}; ERP Simulator & S\&OP planning; purchase orders with goods receipt; order-to-cash pipeline; demand forecasting (SES/DES); invoicing; financial P\&L; ISA-95 B2MML message exchange \\
\bottomrule
\end{tabularx}
\end{table}


\subsection{B2MML Message Integration}

The ISA-95 standard specifies information exchange between Level~3 and Level~4
through structured message types defined in the Business to Manufacturing Markup
Language (B2MML)~\cite{b2mml}.  The SodhiCable prototype simulates these
exchanges through the \tfield{isa95\_messages} table, which logs every
cross-boundary transaction with direction, message type, source system, target
system, payload summary, and processing status.  Table~\ref{tab:b2mml} lists
the B2MML message types implemented in the prototype.

\begin{table}[H]\centering\small
\caption{B2MML Message Types Simulated in the SodhiCable MES}
\label{tab:b2mml}
\begin{tabularx}{\textwidth}{@{} l l l L{5.0cm} @{}}
\toprule
\textbf{Direction} & \textbf{B2MML Message Type} & \textbf{Trigger Event} & \textbf{Payload Description} \\
\midrule
L4$\to$L3 & ProductionSchedule & S\&OP approval; WO release from SO & Approved production plan released to MES for scheduling and dispatch \\
L4$\to$L3 & MaterialRequirement & Manual PO creation & Purchase order for raw materials with quantity, supplier, and expected delivery \\
L3$\to$L4 & MaterialReceipt & Goods receipt processing & Goods receipt confirmation with accepted/rejected quantities and lot number \\
L3$\to$L4 & ProductionPerformance & Shift report submission & OEE metrics, output quantities, scrap totals, and downtime summary per shift \\
% AUTHOR NOTE: Add additional message types as implemented, e.g.:
% L4$\to$L3 & MaterialDefinition & Product/BOM update & Material master data synchronization \\
% L3$\to$L4 & QualityTestResult & Test completion & Spark test, dimensional, and DC resistance results \\
\bottomrule
\end{tabularx}
\end{table}


\subsection{Cross-Level Integration Points}

Three cross-level integration patterns merit particular attention, as they
demonstrate how the MES bridges automation and enterprise layers:

\begin{enumerate}[leftmargin=2em, itemsep=2pt]
\item \textbf{Traceability spanning L1--L4.}  A spark test failure at Level~1
(sensor detects pinhole in \tfield{spark\_test\_log}) propagates upward through
Level~2 (alarm triggered, PLC interlock status set to TRIPPED in the
\texttt{plc\_status} endpoint), Level~3 (NCR auto-created in \tfield{ncr}, lot
quarantined in \tfield{inventory}, forward trace initiated via
\tfield{lot\_tracking}), and Level~4 (yield KPI decremented in
\tfield{shift\_reports}, customer promise date re-evaluated in
\tfield{sales\_orders}).  This end-to-end chain is exercised in the integration
timeline (Section~\ref{sec:timeline}, Events~5--8).

\item \textbf{Recipe--sensor closed loop spanning L2--L3.}  Recipe setpoints
defined at Level~3 (\tfield{recipe\_parameters}: CV zone temperatures, PLCV
pressure, line speed) are pushed to Level~2 as PLC setpoints.  The SCADA PLC
status endpoint (\texttt{/api/scada/plc\_status/\{wc\_id\}}) compares actual
sensor readings against recipe setpoints and computes control error, output
percentage, and interlock status.  Process deviations detected by CUSUM or EWMA
at Level~3 feed back as corrective actions or hold/release decisions.

\item \textbf{AI as intelligence overlay spanning L2--L4.}  The AI/ML layer
(Section~\ref{sec:ai}) operates as a cross-cutting intelligence service
consuming data from all ISA-95 levels: Level~1/2 sensor data for anomaly
detection (z-score, Isolation Forest, Mahalanobis distance), Level~3 production
data for quality prediction (Gradient Boosted Stumps) and maintenance
forecasting (Weibull RUL), and Level~4 demand and financial data for
AI-driven recommendations.  This cross-level positioning distinguishes AI
from the traditional MESA functions, which operate primarily within Level~3.
% AUTHOR NOTE: Consider adding a diagram showing the AI overlay across levels.
\end{enumerate}


%======================================================================
\section{Toward MESA Function~12: An AI/ML Intelligence Layer}\label{sec:mesa12}
%======================================================================

\subsection{Motivation: The MESA-11 Model and the AI Gap}

The MESA International 11-function model, published in 1997~\cite{mesa1997},
defined the canonical functional decomposition of Manufacturing Execution
Systems that has guided MES implementations for nearly three decades.  The
eleven functions---Resource Allocation (F1), Operations Scheduling (F2),
Dispatching (F3), Document Control (F4), Data Collection (F5), Labor Management
(F6), Quality Management (F7), Process Management (F8), Maintenance Management
(F9), Product Tracking (F10), and Performance Analysis (F11)---were conceived in
an era when shop-floor computing consisted of SCADA/DCS systems with limited
storage, batch-mode SPC calculations, and manual analysis of printed control
charts.

The 1997 framework implicitly assumes that human domain experts perform the
analytical synthesis: an engineer reviews control charts (F7), correlates them
with process logs (F8), cross-references maintenance history (F9), and makes
decisions.  The framework provides no function for automated pattern recognition,
predictive modeling, multivariate anomaly detection, or machine-learned
optimization---capabilities that modern AI/ML methods can now deliver at
production-data scale and latency.

MESA International recognized the need for expansion, releasing the Strategic
Initiatives Model in 2004~\cite{mesa2004} and the Smart Manufacturing Model in
subsequent years.  ISA-95 itself was extended with Part~4~\cite{isa95part4},
which added activity models for Manufacturing Operations Management (MOM)
including production operations, quality operations, maintenance operations, and
inventory operations.  However, neither MESA nor ISA-95 has formalized
an AI/ML function within the MES functional hierarchy.

We propose that the original MESA-11 model requires a twelfth function:
\textbf{F12~--- AI/ML Intelligence Layer}, a cross-cutting analytical function
that consumes data produced by F1--F11 and generates predictive, prescriptive,
and diagnostic outputs that augment human decision-making.  The SodhiCable
prototype provides an existence proof through eight implemented AI/ML
capabilities.


\subsection{Implemented AI/ML Capabilities}

The SodhiCable prototype implements eight distinct AI/ML capabilities in pure
Python (no external ML libraries), demonstrating that meaningful intelligence
can be embedded directly in a lightweight MES without requiring a separate data
science platform.  Table~\ref{tab:f12capabilities} summarizes the eight
capabilities and their MESA function touchpoints.

\begin{table}[H]\centering\small
\caption{Proposed F12 Capabilities Implemented in the SodhiCable MES Prototype}
\label{tab:f12capabilities}
\begin{tabularx}{\textwidth}{@{} c L{3.0cm} L{3.5cm} l L{2.5cm} @{}}
\toprule
\textbf{\#} & \textbf{Capability} & \textbf{Method} & \textbf{Touches} & \textbf{Data Source} \\
\midrule
1 & Univariate anomaly detection & Z-score ($|z| > 2.5\sigma$) & F5, F8 & \tfield{process\_data\_live} \\
2 & Multivariate anomaly detection & Isolation Forest (100-tree ensemble) & F5, F8 & \tfield{process\_data\_live} \\
3 & Quality prediction & Gradient Boosted Stumps (50-round additive ensemble) & F7, F11 & \tfield{process\_deviations}, \tfield{scrap\_log}, \tfield{shift\_reports} \\
4 & Multivariate anomaly detection & Mahalanobis distance ($D^2 > \chi^2_{\alpha,p}$) & F5, F7, F8 & \tfield{process\_data\_live} \\
5 & Predictive maintenance & Weibull/exponential RUL estimation & F9 & \tfield{equipment}, \tfield{maintenance}, \tfield{downtime\_log} \\
6 & Statistical process monitoring & CUSUM drift \& EWMA smoothing & F7, F8 & \tfield{process\_deviations}, \tfield{spc\_readings} \\
7 & Prescriptive recommendations & Rules-based engine (7 rule classes) & F1--F11 & Cross-function aggregation \\
8 & Natural language query & Claude API $\to$ SQL generation $\to$ execution & F5, F11 & Full 78-table schema \\
\bottomrule
\end{tabularx}
\end{table}

\subsubsection{Z-Score Anomaly Detection}

The simplest analytical capability: for each (work center, parameter) group in
\tfield{process\_data\_live}, the system computes the sample mean $\bar{x}$ and
standard deviation $s$, then flags any reading where $|z_i| = |x_i - \bar{x}|/s
> 2.5$ as anomalous.  Readings with $|z| > 3.0$ are classified as Critical;
those in the range $2.5 < |z| \leq 3.0$ as Warning.  This extends F5 (data
collection) by adding statistical context to raw sensor values and F8 (process
management) by detecting deviations that may not cross fixed threshold alarms
but represent statistically unusual behavior.

\subsubsection{Isolation Forest}

The prototype implements a pure-Python Isolation Forest~\cite{liu2017}
consisting of 100 isolation trees with a maximum depth of 8 and a subsample
size of 256.  For each (work center, parameter) group, a three-dimensional
feature vector is constructed: $[\,z\text{-score},\;\Delta v / s,\;\Delta v_{-2}
/ s\,]$, where $\Delta v$ is the first difference and $\Delta v_{-2}$ is the
second-order lag difference, both normalized by the parameter's standard
deviation.  The anomaly score is computed as $s(x, n) = 2^{-E[h(x)]/c(n)}$,
where $E[h(x)]$ is the average path length across the 100 trees and $c(n)$ is
the expected path length of an unsuccessful search in a binary search tree of
$n$ items.  Readings with $s > 0.65$ are flagged (Critical if $s > 0.75$).

This method detects interaction effects that univariate z-score analysis
misses---for example, a zone temperature and line speed each within normal range
individually, but whose combination is anomalous.  The implementation requires
approximately 80 lines of Python with no external dependencies.

\subsubsection{Gradient Boosted Decision Stumps}

A 50-round gradient boosting ensemble predicts scrap rate per work center.  Each
round fits a single-split decision stump (one-level tree) to the current
residuals, selecting the feature and split point that minimize mean squared
error.  Four features are extracted per work center: (i)~average deviation
percentage from setpoints, (ii)~count of unresolved process deviations,
(iii)~average OEE, and (iv)~cumulative downtime minutes.  The learning rate is
0.3, and the model trains in-memory on each API call (stateless, reproducible).

Feature importance is computed as the cumulative absolute difference between
left and right split values across all 50 rounds, normalized to percentages.
The RMSE on the training set provides a diagnostic for model fit.  Work centers
with predicted scrap rates above 5\% are flagged as Critical, and those above
3\% as Warning.

\subsubsection{Mahalanobis Distance}

For multivariate anomaly detection that accounts for inter-parameter
correlations, the prototype computes the Mahalanobis distance $D^2$ for each
observation vector:
\[
D^2 = (\mathbf{x} - \boldsymbol{\mu})^\top \, \boldsymbol{\Sigma}^{-1} \, (\mathbf{x} - \boldsymbol{\mu})
\]
where $\boldsymbol{\mu}$ is the sample mean vector and $\boldsymbol{\Sigma}$ is
the sample covariance matrix, computed from all readings for a given work
center.  Observations exceeding the $\chi^2$ critical value at the 99\%
confidence level ($\chi^2_{0.01,2} = 9.21$ for $p=2$ features;
$\chi^2_{0.01,3} = 11.34$ for $p=3$) are flagged as anomalous.

The covariance matrix inversion is implemented analytically for $p = 2$ and
$p = 3$ using the adjugate/determinant method (no NumPy dependency).  For work
centers with more than three tracked parameters, the analysis uses the first
two principal features to maintain tractability in pure Python.

\subsubsection{Weibull/Exponential RUL Prediction}

Remaining Useful Life (RUL) estimation uses a two-model approach.  When
sufficient time-between-failure (TBF) data is available ($\geq 3$
observations), a Weibull distribution is fitted using the method of moments:
\[
F(t) = 1 - e^{-(t/\eta)^\beta}
\]
where the shape parameter $\beta$ is estimated from the coefficient of variation
of the TBF data ($\beta \approx 1/\text{CV}$, clamped to $[\,0.5, \infty)$)
and the scale parameter $\eta$ from the sample mean.  The conditional RUL is
estimated as $E[T - t \mid T > t] \approx E[T] - t$ for the current elapsed
time $t$.  For equipment with insufficient failure history, the model falls back
to the memoryless exponential distribution ($\beta = 1$) with a default
MTBF of 45 days.

The 7-day failure probability $P(\text{fail in } [t, t+168\text{\,hr}])$ is
computed from the Weibull survival function difference, enabling risk-ranked
prioritization of preventive maintenance scheduling.

\subsubsection{CUSUM and EWMA Statistical Monitoring}

Cumulative Sum (CUSUM) and Exponentially Weighted Moving Average (EWMA) charts
extend traditional Shewhart-type SPC (F7/F8) by detecting sustained small
shifts that $\bar{X}$-$R$ charts miss~\cite{montgomery2019}.  The prototype
stores detected CUSUM drifts in \tfield{process\_deviations} with
\tfield{detection\_method = `CUSUM'} and tracks resolution status.  The AI
recommendations engine (Capability~7) monitors unresolved CUSUM drifts and
generates corrective action suggestions referencing likely root causes (die
wear, material variation, calibration drift).

\subsubsection{AI-Driven Recommendations Engine}

A rules-based recommendation engine synthesizes data from across all MESA
functions to generate prioritized, actionable recommendations.  Seven rule
classes are implemented:

\begin{enumerate}[leftmargin=2em, itemsep=1pt]
\item High scrap rate ($>5\%$) at any work center (touches F5, F7, F8, F10)
\item Unresolved CUSUM drift on process parameters (touches F5, F7, F8)
\item Overdue preventive maintenance (touches F1, F9)
\item Low process capability $C_{pk} < 1.33$ (touches F7, F8, F11)
\item Expiring operator certifications within 30 days (touches F1, F6)
\item Declining OEE trend (7-day rolling average) (touches F8, F9, F11)
\item Accumulation of unresolved process deviations ($>3$) (touches F7, F8)
\end{enumerate}

Each recommendation includes a priority level (1 = highest), affected work
center(s), a natural-language explanation with quantitative evidence, and the
MESA function codes it touches.  This cross-function synthesis---combining
quality, maintenance, process, labor, and performance data into a single ranked
action list---is precisely the kind of analytical integration that the original
MESA-11 model left to human judgment.

\subsubsection{Natural Language Query Interface}

The most advanced F12 capability integrates the Anthropic Claude API to provide
a natural language $\to$ SQL $\to$ natural language query pipeline.  The system
constructs a schema-aware prompt containing the full 78-table database structure
and wire-and-cable domain knowledge (OEE targets, $C_{pk}$ thresholds, product
family codes, work center types), then uses Claude's tool-use capability to
iteratively generate and execute SQL queries against the MES database.  Only
\texttt{SELECT} queries are permitted (enforced server-side), and results are
limited to 50 rows.

This capability democratizes MES data access: an operator can ask ``Which lot
should I prioritize?'' and receive a ranked dispatch queue, or a manager can ask
``How much scrap this week?'' and receive an aggregated summary by cause
code---without writing SQL or navigating multiple dashboard screens.  The
agentic loop supports up to 8 iterations of query refinement, enabling complex
multi-step analyses.
% AUTHOR NOTE: Consider adding a concrete query example with input/output
% to illustrate the interaction pattern.


\subsection{F12 as a Cross-Cutting Function}

Unlike functions F1--F11, which each address a distinct operational domain, the
proposed F12 is inherently \textit{cross-cutting}: it consumes data from and
provides analytical outputs to multiple existing functions simultaneously.
Figure~\ref{fig:f12touchpoints} (conceptual; to be drawn) would show F12 at the
center of a hub-and-spoke diagram, with connections to:

\begin{itemize}[leftmargin=2em, itemsep=1pt]
\item \textbf{F5 --- Data Collection:} F12 consumes the sensor data that F5
collects, adding statistical and ML-based anomaly detection as a value-added
analytical layer atop raw data ingestion.
\item \textbf{F7 --- Quality Management:} F12 extends SPC with predictive
quality models (Gradient Boosted Stumps predicting scrap rate from process
features) and multivariate anomaly detection (Mahalanobis distance capturing
inter-parameter correlations that univariate control charts miss).
\item \textbf{F8 --- Process Management:} F12 augments CUSUM/EWMA drift
detection with Isolation Forest anomaly scoring and feeds recommendations back
to process engineers for corrective action.
\item \textbf{F9 --- Maintenance Management:} F12 replaces fixed-interval PM
scheduling with condition-based RUL estimation (Weibull model), enabling
risk-ranked maintenance prioritization.
\item \textbf{F11 --- Performance Analysis:} F12 detects OEE trends and
correlates declining performance with root causes across quality, maintenance,
and process domains---an analytical synthesis that F11's descriptive KPIs alone
cannot provide.
\end{itemize}

This cross-cutting nature is analogous to how ISA-95 Part~4~\cite{isa95part4}
extended the original Parts~1--3 by defining Manufacturing Operations Management
(MOM) activity models that cut across production, quality, maintenance, and
inventory operations.  Just as Part~4 recognized that operational activities do
not exist in functional silos, F12 recognizes that analytical intelligence must
span the full MES data landscape to be effective.
% AUTHOR NOTE: A formal argument for standardization would require mapping F12
% to ISA-95 Part 3 activity models and demonstrating that no existing activity
% adequately covers AI/ML analytics.  Consider drafting this for the conference
% paper.


\subsection{Toward Standardization}

The proposal for MESA Function~12 is not merely academic.  Shojaeinasab et
al.~\cite{shojaeinasab2022} identify five intelligence levels for MES, from
reactive (Level~1) to cognitive (Level~5), and note that most industrial MES
implementations remain at Levels~1--2 (reactive and proactive).  The SodhiCable
prototype operates at Level~3 (predictive) through its Weibull RUL,
Gradient Boosted quality prediction, and anomaly detection capabilities, with
elements of Level~4 (prescriptive) in the recommendations engine.

We suggest that formalizing F12 within the MESA framework would provide three
benefits:

\begin{enumerate}[leftmargin=2em, itemsep=2pt]
\item \textbf{Architectural clarity.}  Vendors and implementers would have a
defined functional boundary for AI/ML capabilities, preventing both
under-investment (treating AI as an afterthought) and scope creep (embedding ML
in every function without coherent architecture).

\item \textbf{Data contract standardization.}  An F12 specification would define
the minimum data interfaces required between the intelligence layer and
functions F1--F11, analogous to how ISA-95 Part~2 defines object models for the
L3--L4 boundary.  This would enable plug-and-play substitution of AI/ML engines
(e.g., replacing the pure-Python Isolation Forest with a scikit-learn or
cloud-based model) without modifying the MES core.

\item \textbf{Maturity assessment.}  An F12 specification could incorporate the
intelligence-level taxonomy of Shojaeinasab et al.~\cite{shojaeinasab2022},
providing a maturity model for organizations to assess and benchmark their MES
AI capabilities.
\end{enumerate}

The SodhiCable prototype demonstrates that a meaningful F12 implementation is
feasible with minimal infrastructure: all eight capabilities run in pure Python
within a single Flask application, with no external ML libraries, GPU
requirements, or separate data platforms.  This accessibility is critical for
the small and mid-size manufacturers (SMMs) that constitute the majority of the
wire and cable industry and for whom heavyweight AI infrastructure is not
economically viable.

% AUTHOR NOTE: For the conference submission, consider:
% (1) Adding empirical validation --- run the AI endpoints against the seeded
%     data and report precision/recall for anomaly detection, RMSE for quality
%     prediction, and RUL calibration accuracy.
% (2) Comparing F12 outputs against human expert decisions on the same data
%     to quantify the value of automated intelligence.
% (3) Surveying MESA International's current working groups to determine if
%     an AI/ML function is already under discussion.
% (4) Discussing computational complexity and scalability --- the pure-Python
%     implementations are O(n * n_trees * max_depth) for Isolation Forest;
%     for production scale, discuss when external ML infrastructure becomes
%     necessary.
 in the main document, or copy sections
% into Eacuello_MES_Design_Prototype.tex at the desired location.
%
% Placeholder citation keys:
%   \cite{mesa1997}      — MESA-11 White Paper #06, 1997
%   \cite{isa95}         — ANSI/ISA-95.00.01-2010 Part 1
%   \cite{isa95part3}    — ANSI/ISA-95.00.03 Part 3 (Activity Models of MOM)
%   \cite{isa95part4}    — ANSI/ISA-95.00.04 Part 4 (Objects and Attributes for MOM)
%   \cite{b2mml}         — B2MML (Business to Manufacturing Markup Language) schema
%   \cite{shojaeinasab2022} — Intelligent MES systematic review
%   \cite{mesa2004}      — MESA Strategic Initiatives Model, 2004
%   \cite{liu2017}       — Isolation Forest (Liu et al., 2008/2012)
%   \cite{montgomery2019} — Montgomery, Introduction to SQC, 8th ed.


%======================================================================
\section{ISA-95 Architecture Mapping}\label{sec:isa95mapping}
%======================================================================

The ISA-95 standard (ANSI/ISA-95.00.01-2010) defines a hierarchical model for
enterprise--control system integration spanning five levels: Level~0 (physical
process), Level~1 (sensing and actuation), Level~2 (supervisory control),
Level~3 (manufacturing operations management), and Level~4 (business planning
and logistics)~\cite{isa95}.  While Section~\ref{sec:overview} introduced this
hierarchy in functional terms, this section provides a formal compliance mapping
between ISA-95 activity domains and the SodhiCable MES implementation modules,
demonstrating how a single Flask-based prototype can simulate the full
automation pyramid.

\subsection{Hierarchical Mapping}

Table~\ref{tab:isa95matrix} presents the ISA-95 compliance matrix, mapping each
level to its activity domain, the corresponding SodhiCable implementation
module(s), and key capabilities realized in the prototype.  Level~3 is further
decomposed by the 11 MESA functions as defined in ISA-95 Part~3's Activity
Models of Manufacturing Operations Management~\cite{isa95part3}.

\begin{table}[H]\centering\small
\caption{ISA-95 Compliance Matrix --- SodhiCable MES Architecture}
\label{tab:isa95matrix}
\begin{tabularx}{\textwidth}{@{} l L{2.4cm} L{3.0cm} L{5.5cm} @{}}
\toprule
\textbf{ISA-95 Level} & \textbf{Activity Domain} & \textbf{Implementation Module(s)} & \textbf{Key Capabilities} \\
\midrule
\multirow{2}{*}{Level 0} & Physical Process & DES Engine & Discrete event simulation of 8-stage factory; compound mixing, wire drawing, extrusion, braiding, cabling, jacketing, testing, packaging \\
\addlinespace[3pt]
\multirow{2}{*}{Level 1} & Sensing \& Actuation & OPC-UA Tag Simulator; \tfield{process\_data\_live}; \tfield{spark\_test\_log}; \tfield{energy\_readings} & 9{,}200+ simulated sensor readings (thermocouples, laser micrometers, pressure transducers, line speed encoders, spark testers); configurable scan interval (default 5\,s) \\
\addlinespace[3pt]
\multirow{2}{*}{Level 2} & Supervisory Control (SCADA/PLC) & \texttt{api\_scada.py}; PID simulation; PLC status endpoint & SCADA drill-down (plant $\to$ work center $\to$ sensor); PLC control-loop visualization (AUTO/MANUAL mode, setpoint tracking, interlock status); alarm management with severity ranking; energy monitoring per work center \\
\addlinespace[3pt]
Level 3 --- F1 & Resource Allocation & \texttt{api\_resources.py} & LP assignment of 26 WCs to WOs; capacity tracking; equipment status management \\
Level 3 --- F2 & Operations Scheduling & \texttt{api\_schedule.py} & SPT/EDD/MILP sequencing; changeover matrix; Gantt visualization \\
Level 3 --- F3 & Dispatching & \texttt{api\_dispatch.py} & Weighted priority dispatch; queue management; real-time job release \\
Level 3 --- F4 & Document Control & \texttt{api\_documents.py} & Recipe versioning (ISA-88 hierarchy); C\,of\,C generation; SOP management \\
Level 3 --- F5 & Data Collection & \texttt{api\_datacollection.py} & Real-time sensor ingestion; threshold alarming; data collection point configuration \\
Level 3 --- F6 & Labor Management & \texttt{api\_labor.py} & Skill matrix; certification tracking; IP shift rostering \\
Level 3 --- F7 & Quality Management & \texttt{api\_quality.py} & $\bar{X}$-$R$ charts; $C_{pk}$ computation; NCR$\to$CAPA workflow; SPC readings \\
Level 3 --- F8 & Process Management & \texttt{api\_process.py} & Recipe execution; CUSUM/EWMA drift detection; PID simulation; hold/release \\
Level 3 --- F9 & Maintenance Mgmt & \texttt{api\_maintenance.py} & MTBF-based PM scheduling; corrective maintenance logging; calibration tracking \\
Level 3 --- F10 & Product Tracking & \texttt{api\_genealogy.py} & Recursive CTE forward/backward trace; lot quarantine; splice zone detection \\
Level 3 --- F11 & Performance Analysis & \texttt{api\_performance.py} & OEE = $A \times P \times Q$; shift reporting; scrap Pareto; KPI dashboards \\
\addlinespace[3pt]
\multirow{2}{*}{Level 4} & Business Planning (ERP) & \texttt{api\_erp.py}; ERP Simulator & S\&OP planning; purchase orders with goods receipt; order-to-cash pipeline; demand forecasting (SES/DES); invoicing; financial P\&L; ISA-95 B2MML message exchange \\
\bottomrule
\end{tabularx}
\end{table}


\subsection{B2MML Message Integration}

The ISA-95 standard specifies information exchange between Level~3 and Level~4
through structured message types defined in the Business to Manufacturing Markup
Language (B2MML)~\cite{b2mml}.  The SodhiCable prototype simulates these
exchanges through the \tfield{isa95\_messages} table, which logs every
cross-boundary transaction with direction, message type, source system, target
system, payload summary, and processing status.  Table~\ref{tab:b2mml} lists
the B2MML message types implemented in the prototype.

\begin{table}[H]\centering\small
\caption{B2MML Message Types Simulated in the SodhiCable MES}
\label{tab:b2mml}
\begin{tabularx}{\textwidth}{@{} l l l L{5.0cm} @{}}
\toprule
\textbf{Direction} & \textbf{B2MML Message Type} & \textbf{Trigger Event} & \textbf{Payload Description} \\
\midrule
L4$\to$L3 & ProductionSchedule & S\&OP approval; WO release from SO & Approved production plan released to MES for scheduling and dispatch \\
L4$\to$L3 & MaterialRequirement & Manual PO creation & Purchase order for raw materials with quantity, supplier, and expected delivery \\
L3$\to$L4 & MaterialReceipt & Goods receipt processing & Goods receipt confirmation with accepted/rejected quantities and lot number \\
L3$\to$L4 & ProductionPerformance & Shift report submission & OEE metrics, output quantities, scrap totals, and downtime summary per shift \\
% AUTHOR NOTE: Add additional message types as implemented, e.g.:
% L4$\to$L3 & MaterialDefinition & Product/BOM update & Material master data synchronization \\
% L3$\to$L4 & QualityTestResult & Test completion & Spark test, dimensional, and DC resistance results \\
\bottomrule
\end{tabularx}
\end{table}


\subsection{Cross-Level Integration Points}

Three cross-level integration patterns merit particular attention, as they
demonstrate how the MES bridges automation and enterprise layers:

\begin{enumerate}[leftmargin=2em, itemsep=2pt]
\item \textbf{Traceability spanning L1--L4.}  A spark test failure at Level~1
(sensor detects pinhole in \tfield{spark\_test\_log}) propagates upward through
Level~2 (alarm triggered, PLC interlock status set to TRIPPED in the
\texttt{plc\_status} endpoint), Level~3 (NCR auto-created in \tfield{ncr}, lot
quarantined in \tfield{inventory}, forward trace initiated via
\tfield{lot\_tracking}), and Level~4 (yield KPI decremented in
\tfield{shift\_reports}, customer promise date re-evaluated in
\tfield{sales\_orders}).  This end-to-end chain is exercised in the integration
timeline (Section~\ref{sec:timeline}, Events~5--8).

\item \textbf{Recipe--sensor closed loop spanning L2--L3.}  Recipe setpoints
defined at Level~3 (\tfield{recipe\_parameters}: CV zone temperatures, PLCV
pressure, line speed) are pushed to Level~2 as PLC setpoints.  The SCADA PLC
status endpoint (\texttt{/api/scada/plc\_status/\{wc\_id\}}) compares actual
sensor readings against recipe setpoints and computes control error, output
percentage, and interlock status.  Process deviations detected by CUSUM or EWMA
at Level~3 feed back as corrective actions or hold/release decisions.

\item \textbf{AI as intelligence overlay spanning L2--L4.}  The AI/ML layer
(Section~\ref{sec:ai}) operates as a cross-cutting intelligence service
consuming data from all ISA-95 levels: Level~1/2 sensor data for anomaly
detection (z-score, Isolation Forest, Mahalanobis distance), Level~3 production
data for quality prediction (Gradient Boosted Stumps) and maintenance
forecasting (Weibull RUL), and Level~4 demand and financial data for
AI-driven recommendations.  This cross-level positioning distinguishes AI
from the traditional MESA functions, which operate primarily within Level~3.
% AUTHOR NOTE: Consider adding a diagram showing the AI overlay across levels.
\end{enumerate}


%======================================================================
\section{Toward MESA Function~12: An AI/ML Intelligence Layer}\label{sec:mesa12}
%======================================================================

\subsection{Motivation: The MESA-11 Model and the AI Gap}

The MESA International 11-function model, published in 1997~\cite{mesa1997},
defined the canonical functional decomposition of Manufacturing Execution
Systems that has guided MES implementations for nearly three decades.  The
eleven functions---Resource Allocation (F1), Operations Scheduling (F2),
Dispatching (F3), Document Control (F4), Data Collection (F5), Labor Management
(F6), Quality Management (F7), Process Management (F8), Maintenance Management
(F9), Product Tracking (F10), and Performance Analysis (F11)---were conceived in
an era when shop-floor computing consisted of SCADA/DCS systems with limited
storage, batch-mode SPC calculations, and manual analysis of printed control
charts.

The 1997 framework implicitly assumes that human domain experts perform the
analytical synthesis: an engineer reviews control charts (F7), correlates them
with process logs (F8), cross-references maintenance history (F9), and makes
decisions.  The framework provides no function for automated pattern recognition,
predictive modeling, multivariate anomaly detection, or machine-learned
optimization---capabilities that modern AI/ML methods can now deliver at
production-data scale and latency.

MESA International recognized the need for expansion, releasing the Strategic
Initiatives Model in 2004~\cite{mesa2004} and the Smart Manufacturing Model in
subsequent years.  ISA-95 itself was extended with Part~4~\cite{isa95part4},
which added activity models for Manufacturing Operations Management (MOM)
including production operations, quality operations, maintenance operations, and
inventory operations.  However, neither MESA nor ISA-95 has formalized
an AI/ML function within the MES functional hierarchy.

We propose that the original MESA-11 model requires a twelfth function:
\textbf{F12~--- AI/ML Intelligence Layer}, a cross-cutting analytical function
that consumes data produced by F1--F11 and generates predictive, prescriptive,
and diagnostic outputs that augment human decision-making.  The SodhiCable
prototype provides an existence proof through eight implemented AI/ML
capabilities.


\subsection{Implemented AI/ML Capabilities}

The SodhiCable prototype implements eight distinct AI/ML capabilities in pure
Python (no external ML libraries), demonstrating that meaningful intelligence
can be embedded directly in a lightweight MES without requiring a separate data
science platform.  Table~\ref{tab:f12capabilities} summarizes the eight
capabilities and their MESA function touchpoints.

\begin{table}[H]\centering\small
\caption{Proposed F12 Capabilities Implemented in the SodhiCable MES Prototype}
\label{tab:f12capabilities}
\begin{tabularx}{\textwidth}{@{} c L{3.0cm} L{3.5cm} l L{2.5cm} @{}}
\toprule
\textbf{\#} & \textbf{Capability} & \textbf{Method} & \textbf{Touches} & \textbf{Data Source} \\
\midrule
1 & Univariate anomaly detection & Z-score ($|z| > 2.5\sigma$) & F5, F8 & \tfield{process\_data\_live} \\
2 & Multivariate anomaly detection & Isolation Forest (100-tree ensemble) & F5, F8 & \tfield{process\_data\_live} \\
3 & Quality prediction & Gradient Boosted Stumps (50-round additive ensemble) & F7, F11 & \tfield{process\_deviations}, \tfield{scrap\_log}, \tfield{shift\_reports} \\
4 & Multivariate anomaly detection & Mahalanobis distance ($D^2 > \chi^2_{\alpha,p}$) & F5, F7, F8 & \tfield{process\_data\_live} \\
5 & Predictive maintenance & Weibull/exponential RUL estimation & F9 & \tfield{equipment}, \tfield{maintenance}, \tfield{downtime\_log} \\
6 & Statistical process monitoring & CUSUM drift \& EWMA smoothing & F7, F8 & \tfield{process\_deviations}, \tfield{spc\_readings} \\
7 & Prescriptive recommendations & Rules-based engine (7 rule classes) & F1--F11 & Cross-function aggregation \\
8 & Natural language query & Claude API $\to$ SQL generation $\to$ execution & F5, F11 & Full 78-table schema \\
\bottomrule
\end{tabularx}
\end{table}

\subsubsection{Z-Score Anomaly Detection}

The simplest analytical capability: for each (work center, parameter) group in
\tfield{process\_data\_live}, the system computes the sample mean $\bar{x}$ and
standard deviation $s$, then flags any reading where $|z_i| = |x_i - \bar{x}|/s
> 2.5$ as anomalous.  Readings with $|z| > 3.0$ are classified as Critical;
those in the range $2.5 < |z| \leq 3.0$ as Warning.  This extends F5 (data
collection) by adding statistical context to raw sensor values and F8 (process
management) by detecting deviations that may not cross fixed threshold alarms
but represent statistically unusual behavior.

\subsubsection{Isolation Forest}

The prototype implements a pure-Python Isolation Forest~\cite{liu2017}
consisting of 100 isolation trees with a maximum depth of 8 and a subsample
size of 256.  For each (work center, parameter) group, a three-dimensional
feature vector is constructed: $[\,z\text{-score},\;\Delta v / s,\;\Delta v_{-2}
/ s\,]$, where $\Delta v$ is the first difference and $\Delta v_{-2}$ is the
second-order lag difference, both normalized by the parameter's standard
deviation.  The anomaly score is computed as $s(x, n) = 2^{-E[h(x)]/c(n)}$,
where $E[h(x)]$ is the average path length across the 100 trees and $c(n)$ is
the expected path length of an unsuccessful search in a binary search tree of
$n$ items.  Readings with $s > 0.65$ are flagged (Critical if $s > 0.75$).

This method detects interaction effects that univariate z-score analysis
misses---for example, a zone temperature and line speed each within normal range
individually, but whose combination is anomalous.  The implementation requires
approximately 80 lines of Python with no external dependencies.

\subsubsection{Gradient Boosted Decision Stumps}

A 50-round gradient boosting ensemble predicts scrap rate per work center.  Each
round fits a single-split decision stump (one-level tree) to the current
residuals, selecting the feature and split point that minimize mean squared
error.  Four features are extracted per work center: (i)~average deviation
percentage from setpoints, (ii)~count of unresolved process deviations,
(iii)~average OEE, and (iv)~cumulative downtime minutes.  The learning rate is
0.3, and the model trains in-memory on each API call (stateless, reproducible).

Feature importance is computed as the cumulative absolute difference between
left and right split values across all 50 rounds, normalized to percentages.
The RMSE on the training set provides a diagnostic for model fit.  Work centers
with predicted scrap rates above 5\% are flagged as Critical, and those above
3\% as Warning.

\subsubsection{Mahalanobis Distance}

For multivariate anomaly detection that accounts for inter-parameter
correlations, the prototype computes the Mahalanobis distance $D^2$ for each
observation vector:
\[
D^2 = (\mathbf{x} - \boldsymbol{\mu})^\top \, \boldsymbol{\Sigma}^{-1} \, (\mathbf{x} - \boldsymbol{\mu})
\]
where $\boldsymbol{\mu}$ is the sample mean vector and $\boldsymbol{\Sigma}$ is
the sample covariance matrix, computed from all readings for a given work
center.  Observations exceeding the $\chi^2$ critical value at the 99\%
confidence level ($\chi^2_{0.01,2} = 9.21$ for $p=2$ features;
$\chi^2_{0.01,3} = 11.34$ for $p=3$) are flagged as anomalous.

The covariance matrix inversion is implemented analytically for $p = 2$ and
$p = 3$ using the adjugate/determinant method (no NumPy dependency).  For work
centers with more than three tracked parameters, the analysis uses the first
two principal features to maintain tractability in pure Python.

\subsubsection{Weibull/Exponential RUL Prediction}

Remaining Useful Life (RUL) estimation uses a two-model approach.  When
sufficient time-between-failure (TBF) data is available ($\geq 3$
observations), a Weibull distribution is fitted using the method of moments:
\[
F(t) = 1 - e^{-(t/\eta)^\beta}
\]
where the shape parameter $\beta$ is estimated from the coefficient of variation
of the TBF data ($\beta \approx 1/\text{CV}$, clamped to $[\,0.5, \infty)$)
and the scale parameter $\eta$ from the sample mean.  The conditional RUL is
estimated as $E[T - t \mid T > t] \approx E[T] - t$ for the current elapsed
time $t$.  For equipment with insufficient failure history, the model falls back
to the memoryless exponential distribution ($\beta = 1$) with a default
MTBF of 45 days.

The 7-day failure probability $P(\text{fail in } [t, t+168\text{\,hr}])$ is
computed from the Weibull survival function difference, enabling risk-ranked
prioritization of preventive maintenance scheduling.

\subsubsection{CUSUM and EWMA Statistical Monitoring}

Cumulative Sum (CUSUM) and Exponentially Weighted Moving Average (EWMA) charts
extend traditional Shewhart-type SPC (F7/F8) by detecting sustained small
shifts that $\bar{X}$-$R$ charts miss~\cite{montgomery2019}.  The prototype
stores detected CUSUM drifts in \tfield{process\_deviations} with
\tfield{detection\_method = `CUSUM'} and tracks resolution status.  The AI
recommendations engine (Capability~7) monitors unresolved CUSUM drifts and
generates corrective action suggestions referencing likely root causes (die
wear, material variation, calibration drift).

\subsubsection{AI-Driven Recommendations Engine}

A rules-based recommendation engine synthesizes data from across all MESA
functions to generate prioritized, actionable recommendations.  Seven rule
classes are implemented:

\begin{enumerate}[leftmargin=2em, itemsep=1pt]
\item High scrap rate ($>5\%$) at any work center (touches F5, F7, F8, F10)
\item Unresolved CUSUM drift on process parameters (touches F5, F7, F8)
\item Overdue preventive maintenance (touches F1, F9)
\item Low process capability $C_{pk} < 1.33$ (touches F7, F8, F11)
\item Expiring operator certifications within 30 days (touches F1, F6)
\item Declining OEE trend (7-day rolling average) (touches F8, F9, F11)
\item Accumulation of unresolved process deviations ($>3$) (touches F7, F8)
\end{enumerate}

Each recommendation includes a priority level (1 = highest), affected work
center(s), a natural-language explanation with quantitative evidence, and the
MESA function codes it touches.  This cross-function synthesis---combining
quality, maintenance, process, labor, and performance data into a single ranked
action list---is precisely the kind of analytical integration that the original
MESA-11 model left to human judgment.

\subsubsection{Natural Language Query Interface}

The most advanced F12 capability integrates the Anthropic Claude API to provide
a natural language $\to$ SQL $\to$ natural language query pipeline.  The system
constructs a schema-aware prompt containing the full 78-table database structure
and wire-and-cable domain knowledge (OEE targets, $C_{pk}$ thresholds, product
family codes, work center types), then uses Claude's tool-use capability to
iteratively generate and execute SQL queries against the MES database.  Only
\texttt{SELECT} queries are permitted (enforced server-side), and results are
limited to 50 rows.

This capability democratizes MES data access: an operator can ask ``Which lot
should I prioritize?'' and receive a ranked dispatch queue, or a manager can ask
``How much scrap this week?'' and receive an aggregated summary by cause
code---without writing SQL or navigating multiple dashboard screens.  The
agentic loop supports up to 8 iterations of query refinement, enabling complex
multi-step analyses.
% AUTHOR NOTE: Consider adding a concrete query example with input/output
% to illustrate the interaction pattern.


\subsection{F12 as a Cross-Cutting Function}

Unlike functions F1--F11, which each address a distinct operational domain, the
proposed F12 is inherently \textit{cross-cutting}: it consumes data from and
provides analytical outputs to multiple existing functions simultaneously.
Figure~\ref{fig:f12touchpoints} (conceptual; to be drawn) would show F12 at the
center of a hub-and-spoke diagram, with connections to:

\begin{itemize}[leftmargin=2em, itemsep=1pt]
\item \textbf{F5 --- Data Collection:} F12 consumes the sensor data that F5
collects, adding statistical and ML-based anomaly detection as a value-added
analytical layer atop raw data ingestion.
\item \textbf{F7 --- Quality Management:} F12 extends SPC with predictive
quality models (Gradient Boosted Stumps predicting scrap rate from process
features) and multivariate anomaly detection (Mahalanobis distance capturing
inter-parameter correlations that univariate control charts miss).
\item \textbf{F8 --- Process Management:} F12 augments CUSUM/EWMA drift
detection with Isolation Forest anomaly scoring and feeds recommendations back
to process engineers for corrective action.
\item \textbf{F9 --- Maintenance Management:} F12 replaces fixed-interval PM
scheduling with condition-based RUL estimation (Weibull model), enabling
risk-ranked maintenance prioritization.
\item \textbf{F11 --- Performance Analysis:} F12 detects OEE trends and
correlates declining performance with root causes across quality, maintenance,
and process domains---an analytical synthesis that F11's descriptive KPIs alone
cannot provide.
\end{itemize}

This cross-cutting nature is analogous to how ISA-95 Part~4~\cite{isa95part4}
extended the original Parts~1--3 by defining Manufacturing Operations Management
(MOM) activity models that cut across production, quality, maintenance, and
inventory operations.  Just as Part~4 recognized that operational activities do
not exist in functional silos, F12 recognizes that analytical intelligence must
span the full MES data landscape to be effective.
% AUTHOR NOTE: A formal argument for standardization would require mapping F12
% to ISA-95 Part 3 activity models and demonstrating that no existing activity
% adequately covers AI/ML analytics.  Consider drafting this for the conference
% paper.


\subsection{Toward Standardization}

The proposal for MESA Function~12 is not merely academic.  Shojaeinasab et
al.~\cite{shojaeinasab2022} identify five intelligence levels for MES, from
reactive (Level~1) to cognitive (Level~5), and note that most industrial MES
implementations remain at Levels~1--2 (reactive and proactive).  The SodhiCable
prototype operates at Level~3 (predictive) through its Weibull RUL,
Gradient Boosted quality prediction, and anomaly detection capabilities, with
elements of Level~4 (prescriptive) in the recommendations engine.

We suggest that formalizing F12 within the MESA framework would provide three
benefits:

\begin{enumerate}[leftmargin=2em, itemsep=2pt]
\item \textbf{Architectural clarity.}  Vendors and implementers would have a
defined functional boundary for AI/ML capabilities, preventing both
under-investment (treating AI as an afterthought) and scope creep (embedding ML
in every function without coherent architecture).

\item \textbf{Data contract standardization.}  An F12 specification would define
the minimum data interfaces required between the intelligence layer and
functions F1--F11, analogous to how ISA-95 Part~2 defines object models for the
L3--L4 boundary.  This would enable plug-and-play substitution of AI/ML engines
(e.g., replacing the pure-Python Isolation Forest with a scikit-learn or
cloud-based model) without modifying the MES core.

\item \textbf{Maturity assessment.}  An F12 specification could incorporate the
intelligence-level taxonomy of Shojaeinasab et al.~\cite{shojaeinasab2022},
providing a maturity model for organizations to assess and benchmark their MES
AI capabilities.
\end{enumerate}

The SodhiCable prototype demonstrates that a meaningful F12 implementation is
feasible with minimal infrastructure: all eight capabilities run in pure Python
within a single Flask application, with no external ML libraries, GPU
requirements, or separate data platforms.  This accessibility is critical for
the small and mid-size manufacturers (SMMs) that constitute the majority of the
wire and cable industry and for whom heavyweight AI infrastructure is not
economically viable.

% AUTHOR NOTE: For the conference submission, consider:
% (1) Adding empirical validation --- run the AI endpoints against the seeded
%     data and report precision/recall for anomaly detection, RMSE for quality
%     prediction, and RUL calibration accuracy.
% (2) Comparing F12 outputs against human expert decisions on the same data
%     to quantify the value of automated intelligence.
% (3) Surveying MESA International's current working groups to determine if
%     an AI/ML function is already under discussion.
% (4) Discussing computational complexity and scalability --- the pure-Python
%     implementations are O(n * n_trees * max_depth) for Isolation Forest;
%     for production scale, discuss when external ML infrastructure becomes
%     necessary.
